\documentclass[11pt,letterpaper]{article}
\usepackage[margin=0.88in,headheight=14pt]{geometry}
\usepackage[T1]{fontenc}
\usepackage{lmodern,microtype}
\usepackage{amsmath,amssymb,amsthm}
\usepackage{booktabs,longtable,array,tabularx}
\usepackage{enumitem,xcolor,fancyhdr,listings,needspace}
\usepackage[hyphens]{url}
\usepackage{hyperref}
\definecolor{ink}{HTML}{193747}
\definecolor{linkblue}{HTML}{155E83}
\hypersetup{colorlinks=true,linkcolor=ink,urlcolor=linkblue,citecolor=ink,
  pdftitle={A Smaller Validated Checker for the Concise iPRA Axiomatization},
  pdfauthor={Prepared for Vladimir Reshetnikov}}
\setlength{\parskip}{0.35em}
\setlength{\parindent}{0pt}
\setlist{nosep,topsep=0.35em,leftmargin=*}
\emergencystretch=2em
\newtheorem{theorem}{Theorem}[section]
\newtheorem{lemma}[theorem]{Lemma}
\newtheorem{proposition}[theorem]{Proposition}
\newtheorem{corollary}[theorem]{Corollary}
\theoremstyle{definition}
\newtheorem{definition}[theorem]{Definition}
\theoremstyle{remark}
\newtheorem{remark}[theorem]{Remark}
\newcommand{\D}{\mathsf D}
\newcommand{\Q}{\mathsf Q}
\newcommand{\ipra}{\mathrm{iPRA}}
\pagestyle{fancy}
\fancyhf{}
\fancyhead[L]{\small Concise iPRA}
\fancyhead[R]{\small A smaller validated checker}
\fancyfoot[C]{\thepage}
\lstset{language=Python,basicstyle=\ttfamily\scriptsize,breaklines=true,
  columns=fullflexible,keepspaces=true,showstringspaces=false,
  frame=single,rulecolor=\color{ink!25},framesep=5pt}

\title{\color{ink}A Smaller Validated Checker\\for the Concise iPRA Axiomatization}
\author{Research note prepared for Vladimir Reshetnikov}
\date{11 November 2026}
\begin{document}
\maketitle
\begin{abstract}
The verified improvement in this note is a smaller complete implementation
of the existing ConciseIPRA presentation. The fully validated Eq/Q/D checker
has 870 normalized Python bytes, 383 significant tokens, and 442 AST nodes,
against 897 bytes, 401 tokens, and 468 nodes in the newest matching
repository archive. Its maximum-compression gzip size is 507 bytes, compared
with 505. The Eq/D public-language variant has 849 bytes, compared with 878.
The reduction replaces the recursive raw-type mode of the validator by an
exact characterization of integer/list JSON trees using their Python
representation. The arithmetic proof kernel is unchanged.

This is a modest reduction under the complete-proof-checker criterion.
It is not a new smaller set of arithmetic axioms, a mathematical breakthrough,
an absolute minimum, or an improvement over the separate 521-byte guarded
checker, which has a different validation contract. The note states the
complete inherited axiomatization, both interpretations and their composites,
the new validator-preservation proof, measured comparisons, and reproducible
checks. Attributed proof appendices make the mathematical dependencies
available inside the archive.
\end{abstract}

\section*{Result at a glance}
\begin{center}
\begin{tabular}{lrrrr}
\toprule
Full Eq/Q/D presentation & Bytes & Tokens & AST & Gzip-9\\
\midrule
Pinned incoming control & 897 & 401 & 468 & 505\\
This presentation & \textbf{870} & \textbf{383} & \textbf{442} & 507\\
Reduction & 27 & 18 & 26 & $-2$\\
\bottomrule
\end{tabular}
\end{center}
All four columns count a complete standalone program, including parsing,
validation, proof checking, and the acceptance output. Gzip is informational.
The negative last entry means that compressed size increases by two bytes.

\clearpage
\tableofcontents
\clearpage

\section{The question and the comparison boundary}
\label{sec:scope}
The requested objective is a simpler axiomatization of full first-order
intuitionistic primitive recursive arithmetic, with equiconsistency and
bi-interpretability, while allowing a different language or numerical
domain. The recurring operational simplicity criterion is the size of a
complete, straightforwardly inspectable proof checker. We distinguish that
criterion from the number of primitive mathematical rules and from the
size of a proof certificate.

The comparison is pinned to ProveIt commit
\texttt{42a48d2f04ad72036e8d7e365ea91ea08c76a9e3}, specifically
\path{Logic/PrimitiveRecursiveArithmetic/ConciseIPRA}.
The current README describes 38 research packages and a Lean development.
In addition, the incoming directory contains the newer
\path{ipra_one_sided_transport_2026-11-11.zip} and
\path{ipra_validated_core_2026-11-11.zip}. Both were inspected before
selecting the control. The latter supplies the 897-byte full native and
878-byte Eq/D checkers used here. Exact repository identifiers, copied
source hashes, and the incoming Git blob identities are recorded in
\path{reference/provenance.json}.

The full native public language has equality, Q, and D. The typed input
contract also requires every node of the whole certificate, including an
unused prefix, to be an exact Python integer or list after JSON parsing.
In particular, Boolean values, floating-point values, strings, dictionaries,
and null are rejected. The advertised goal must be a canonical closed
formula. A shorter program obtained by omitting one of these checks would
not be a matching improvement.

The separate compact guarded envelope is 256 significant tokens, 278 AST
nodes, and 521 normalized bytes, attained by three different programs.
Those programs assume a canonical guarded public goal. None of the present
claims compares 870 against 521 as though the input contracts were equal.
The current seven-rule mathematical variants also remain available; this
note does not assert that its unchanged mathematical basis has fewer rules.

\paragraph{Status of the requested mathematical simplification.}
No further reduction of the arithmetic or logical axiom basis was established
in this investigation. Removing the initializer, eliminating primitive
successor, replacing case introduction by uniform progress, and narrowing
transport have already been extensively investigated in the pinned sources.
The positive result below is a proved simplification of a complete validated
presentation of the same theory. This distinction is part of the result,
not an implicit claim of a new foundation.

\section{The mathematical theory implemented by the checker}
\label{sec:unchanged-native-theory}

The reduction in this report concerns the complete validated proof checker.
It does not remove a mathematical axiom, weaken a side condition, or introduce
a new arithmetic principle. For clarity, this section states the existing
theory in full and identifies the inherited proof of its relationship with
iPRA. The source is the ConciseIPRA snapshot at commit
\texttt{42a48d2f04ad72036e8d7e365ea91ea08c76a9e3}.
The attributed foundational proofs are reproduced in the appendices.

The presentation below is the eight-rule calculus called \texttt{IPrf} in
the repository. The existing \texttt{NPrf} kernel also has eight rules, with
weaker equational transport and direct equation-progress introduction;
\path{iprf_iff_nprf} proves that their consequence relations coincide.
The repository additionally has equivalent seven-rule presentations.
None of those mathematical rule counts is improved here.

\subsection{Finite syntax and all primitive rules}

There is one numerical sort. Its terms are
\[
 t ::= x\mid 0\mid S(t)\mid I(n,a;u.h).
\]
Only $u$ is bound in the update body $h$; it is not bound in $n$ or $a$.
Every other parameter of $h$ stays fixed during the iteration. A body is a
fixed finite expression, not a quantified code for a function. In particular,
this syntax supplies no uniform evaluator for arbitrary primitive-recursive
descriptions. All substitutions below avoid capture, and bound variables may
be renamed.

The native formulas are
\[
 A ::= (s=t)\mid Q_yA\mid D_{x,y}A.
\]
The displayed variables are bound in the matrix of their respective binder;
$x$ and $y$ in $D$ are distinct. A context $\Gamma$ is a finite reusable list
of formulas. Write $1:=S0$ and designate $0=1$ as native contradiction.
Neither ordinary implication nor ordinary falsity is primitive syntax here.

Here is the entire primitive inventory. The context in each conclusion is
the same retained context as in its premises.
\begin{align*}
 &\frac{}{\Gamma\vdash A}\quad(A\in\Gamma)
 &&\mathrm{Hyp},\\[4pt]
 &\frac{}{\Gamma\vdash I(0,a;u.h)=a}
 &&I_0,\\[4pt]
 &\frac{}{\Gamma\vdash I(Sn,a;u.h)=h[I(n,a;u.h)/u]}
 &&I_S,\\[4pt]
 &\frac{\Gamma\vdash s=t\qquad\Gamma\vdash F(s)}
        {\Gamma\vdash F(t)}
 &&\mathrm{Tr},\\[4pt]
 &\frac{\Gamma\vdash Q_yB(y)\qquad\Gamma\vdash B(0)}
        {\Gamma\vdash B(t)}
 &&\mathrm{QE},\\[4pt]
 &\frac{\Gamma\vdash D_{x,y}B(x,y)\qquad
         \Gamma,Q_yB(b,y)\vdash A}
        {\Gamma\vdash A}
 &&\mathrm{DE},\\[4pt]
 &\frac{\Gamma,B(t,0)\vdash B(t,Sz)}
        {\Gamma\vdash D_{x,y}B(x,y)}
 &&\mathrm{DI}_{\mathrm{case}},\\[4pt]
 &\frac{\Gamma,E(t,z)\vdash E(t,Sz)}
        {\Gamma\vdash D_{x,y}E(x,y)}
 &&\mathrm{DI}_{\mathrm{progress}}.
\end{align*}
In the last rule $E$ must be a \emph{literal equation}: its outermost
constructor is equality. An equation occurring inside $Q$ or $D$ does not
meet that condition. In Tr, $F(v)$ is any native formula template with one
designated free numerical variable; that variable may occur any finite number
of times, including zero. It is renamed away from the substituted terms and
the retained context. The term $t$ in QE is arbitrary.

In DE, the eigenvariable $b$ is fresh for $\Gamma$, the conclusion $A$,
and the passive parameters of $B$. In either DI rule, $z$ is fresh for
$\Gamma$, the witness $t$, and the passive parameters of the matrix.
The displayed local assumption is discharged. These freshness conditions
prevent a witness from depending on an eigenvariable introduced later.

There are no other primitive equality axioms, successor axioms, induction
schemas, logical connectives, or logical rules. Weakening, contraction,
exchange, fresh renaming, substitution, and assumption cut are derived
operations. For example, if $r=I(0,t;u.0)$, then $I_0$ proves $r=t$;
using that equation as both premises of Tr with $F(v)\equiv(v=t)$ proves
$t=t$. Symmetry, transitivity, and congruence then follow by the usual
equality templates. The full derivations are in
Section~\ref{sec:choice-introduction}.

\subsection{Where the ordinary logical operations come from}

The scalar selector
\[
 C(v,a,b):=I(v,a;u.b),\qquad
 p(y):=I(y,0;u.C(u,1,0))
\]
has $C(0,a,b)=a$ and $C(Sv,a,b)=b$; the update of $C$ ignores $u$.
The term $p$ is the parity iterator, and only its fixed instances
$p(1)=1$ and $p(2)=0$ are needed for the conjunction macro below.

A finite syntactic gate $G_vA$ and selector $M_v(A,B)$ satisfy
\[
 \vdash G_0A,\qquad G_{Sv}A\dashv\vdash A,
 \qquad M_0(A,B)\dashv\vdash A,
 \qquad M_{Sv}(A,B)\dashv\vdash B.
\]
Here interderivability holds under arbitrary additional retained assumptions.
At equations their definitions are
\[
 G_v(s=t):=(C(v,0,s)=C(v,0,t)),
\]
\[
 M_v(s=t,r=q):=(C(v,s,r)=C(v,t,q)).
\]
For a completely finite extension to all formulas, first replace each
$Q_yB$ by the interderivable $D_{r,y}B$, with $r$ unused and fresh.
The gate commutes with $D$. To align selectors of unequal binder depth,
pad an equation $E$ by $D_{r,y}G_yE$; then recurse under the aligned $D$
binders. Section~\ref{subsec:choice-logic} gives the exact terminating
recursion, the padding equivalence, and its proof. This is external recursion
on fixed formula syntax, not an internal truth predicate.

Using those operations, one may take the ordinary logical macros to be
\[
\begin{aligned}
 \bot_*&:=(0=1),\\
 A\mathbin{\to_*}B&:=Q_yM_y(A,B),\\
 A\mathbin{\wedge_*}B&:=Q_yG_yM_{p(y)}(A,B),\\
 \exists_*x\,A(x)&:=D_{x,y}G_yA(x),\\
 A\mathbin{\vee_*}B&:=\exists_*v\,M_v(A,B),\\
 \forall_*x\,A(x)&:=A(0)\mathbin{\wedge_*}Q_xA(x).
\end{aligned}
\]
All new bound variables are fresh; before taking a selector or gate, use
the equivalent $D$-only version just described when necessary.
The equivalent $D$-only inventory in Section~\ref{subsec:choice-logic}
also avoids treating $Q$ as a public formula constructor.

These macros have every intuitionistic natural-deduction rule. Implication
introduction uses the zero selector as its local antecedent and supplies the
successor selector as the consequent; elimination specializes a closure at
$1$. Existential introduction uses its supplied witness, while DE opens a
fresh witness for existential elimination; the zero gate supplies the required
base. Conjunction projections use closure instances $1$ and $2$. Disjunction
uses witnesses $0$ and $1$, followed by the derived numerical case rule for
elimination. A fresh-parameter proof gives universal introduction, and its
zero instance together with its closure gives arbitrary specialization.
Under $0=1$, the selector equations give
$s=C(0,s,t)=C(1,s,t)=t$; external formula induction then proves explosion.
No excluded middle for quantified formulas is introduced.

The resulting numerical case principle and literal-equation induction are
\[
 \frac{\Gamma\vdash A(0)\qquad\Gamma\vdash A(Sz)}
      {\Gamma\vdash A(t)},
 \qquad
 \frac{\Gamma\vdash E(0)\qquad\Gamma,E(z)\vdash E(Sz)}
      {\Gamma\vdash E(t)},
\]
with the usual fresh $z$ and with $E$ a literal equation. In the second
construction, progress DI with an unused witness supplies a closure, and QE
applies it to the given base. The first construction uses case DI and QE.
One must not silently apply the first rule with $z$ free in a retained
assumption. The inherited proofs explicitly discharge that assumption using
the implication macro or use coverage with equality hypotheses instead.

\subsection{The ordinary first-order realization}

To state bi-interpretability in ordinary first-order terms, elaborate each
fixed iterator body $h(u,\vec p)$ to a staged function symbol
$i_h(n,a,\vec p)$ satisfying
\[
 i_h(0,a,\vec p)=a,\qquad
 i_h(Sn,a,\vec p)=h(i_h(n,a,\vec p),\vec p).
\]
Only finitely described bodies are introduced at each stage. This elaboration
changes notation, not the numerical domain.

Let $T_{\mathrm{eq}}$ be intuitionistic first-order logic with actual equality
in this signature, the displayed computation schemas, and
\[
 \forall x\,(x=0\vee\exists y\,x=Sy),
 \qquad (0=S0)\to\bot.
\]
It also has induction for every literal equation $E(n,\vec p)$:
\[
 \bigl(E(0,\vec p)\wedge
       \forall n\,(E(n,\vec p)\to E(Sn,\vec p))\bigr)
       \to\forall n\,E(n,\vec p),
\]
universally closed over the passive parameters. Its contextual induction
rule is the equivalent displayed rule above. The separation sentence is
essential: it connects native contradiction to genuine ordinary falsity.
The equation-only part by itself must not be substituted for this theory.

The ordinary reading $\rho$ fixes terms and equality and is given by
\[
 \rho(Q_yB)=\rho(B(0))\to\forall y\,\rho(B(y)),
\]
\[
 \rho(D_{x,y}B)=
   \exists x\bigl(\rho(B(x,0))\to\forall y\,\rho(B(x,y))\bigr).
\]
Case DI is sound by coverage. Progress DI is sound by equation induction
under a temporarily retained base equation. The other rules have their
ordinary equality or logical meanings. Conversely, the logical macros,
derived numerical cases, and equation induction recover the ordinary
realization. The native translation $\lambda$ of ordinary formulas fixes
equations and replaces the six ordinary logical operations by the macros
above. Sections~\ref{sec:choice-introduction} and~\ref{sec:ordinary-bridge}
prove preservation of sequents and both formula composites, not merely
agreement in the standard natural numbers.

\subsection{The arithmetic bridge to standard iPRA}

Here standard iPRA means intuitionistic first-order arithmetic with one
function symbol for each fixed primitive-recursive description: zero,
successor, projections, composition, and primitive recursion. It has their
defining equations, $Sx\ne0$, successor injectivity, and induction for every
quantifier-free formula, with parameters. The usual intuitionistic logical
rules retain assumptions and enforce fresh variables for universal
introduction and existential elimination.

The inherited arithmetic bridge has a necessary dependency order.
First, a scalar term $P$ is constructed and proved to satisfy
$P0=0$ and $P(Sn)=n$, using no pre-existing pairing or equality detector
(Section~\ref{sec:hybrid-bootstrap}). Next, truncated subtraction and a
Boolean term reflecting actual numerical equality are constructed.
Every fixed quantifier-free formula is then equivalent to its Boolean
characteristic equation, so equation induction yields quantifier-free
induction (Section~\ref{sec:eic-completion}). Only after that step are pairing
and its two projections used to compile general primitive recursion
(Section~\ref{sec:iterator-coding}).

For example, let $\pi,L,H$ satisfy
$L(\pi(a,b))=a$ and $H(\pi(a,b))=b$. A target recursion
\[
 f(0,\vec p)=g(\vec p),\qquad
 f(Sn,\vec p)=h(n,f(n,\vec p),\vec p)
\]
is compiled, after its earlier descriptions $g,h$ have been compiled, by
\[
 \begin{aligned}
 W(u,\vec p)&:=
   \pi\bigl(SL(u),h^*(L(u),H(u),\vec p)\bigr),\\
 V(n,\vec p)&:=I(n,\pi(0,g^*(\vec p));u.W(u,\vec p)),\\
 f^*(n,\vec p)&:=H(V(n,\vec p)).
 \end{aligned}
\]
Equation induction proves $L(V(n,\vec p))=n$, after which the projection
equations prove both target defining equations. The compiler is external
recursion on each finite description. Pairing stores a finite state; the
interpretation domain is not restricted to the image of $\pi$.

In the other direction, iPRA expresses every scalar iterator by ordinary
primitive recursion whose update ignores the recursion counter. It proves
coverage, separation, and equation induction. Thus each theory interprets
the other, and the following stronger statement supplies the required
round trips.

\begin{theorem}[Inherited whole-domain bi-interpretation]
\label{thm:inherited-biinterpretation}
There are effective translations $\Phi$ from native formulas to iPRA and
$\Psi$ from iPRA to native formulas such that, for all finite contexts,
\[
 \Gamma\vdash A
 \quad\Longleftrightarrow\quad
 \Phi\Gamma\vdash_{\mathrm{iPRA}}\Phi A,
\]
\[
 \Delta\vdash_{\mathrm{iPRA}}\varphi
 \quad\Longleftrightarrow\quad
 \Psi\Delta\vdash\Psi\varphi.
\]
Every native $A$ is interderivable with $\Psi\Phi A$, and iPRA proves
$\varphi\leftrightarrow\Phi\Psi\varphi$ for every ordinary formula
$\varphi$. Both interpretations use the entire numerical domain and actual
equality. In each direction the internally definable numerical return
isomorphism has graph $x=y$. In particular,
\[
 \varnothing\vdash 0=S0
 \quad\Longleftrightarrow\quad
 \varnothing\vdash_{\mathrm{iPRA}}\bot.
\]
\end{theorem}

\begin{proof}
Compose the native/ordinary translations $\rho,\lambda$ with the two
arithmetic compilations just described. For a fixed iterator, equation
induction proves uniqueness from its seed and update equation. The same
argument proves uniqueness of a fixed ordinary primitive recursion.
External induction on a finite description therefore shows that its double
translation equals its original value at every internal argument. This is
an internal numerical equality, valid at nonstandard arguments as well as
standard numerals.

The identity numerical graph is consequently total, bijective, and preserves
every interpreted operation. The logical macro equivalences and binder
composites then give both formula round trips by external structural
induction. In particular, $D_{x,y}B$ is interderivable with the native
existential of the native implication from $B(x,0)$ to the native universal
of $B(x,y)$; the analogous closure statement holds for $Q$. The explicit
ordinary separation sentence and native explosion preserve contradiction.
Finite proof translations establish equiconsistency. The detailed
constructions are in Sections~\ref{sec:choice-introduction},
\ref{sec:ordinary-bridge}, and~\ref{sec:iterator-coding}.
\end{proof}

The pinned Lean development records the ordinary arithmetic bridge as
\path{biInterp}, \path{teq_roundTrip}, and
\path{pra_roundTrip}; the native bridge is recorded by
\path{native_iff_pra}, \path{pra_iff_native},
\path{native_roundTrip}, \path{pra_native_roundTrip}, and
\path{native_inconsistent_iff} in
\path{NativeBiInterpretation.lean} and its dependencies.
The comparison \path{iprf_iff_nprf} connects the eight rules stated here
to that kernel. These are attributed existing results; the present
validation refinement neither adds a new Lean theorem nor claims that the
Python implementation has been formally verified in Lean. Its separate
checking argument concerns the representation and validation of certificates
for this unchanged mathematics.


\section{A smaller validation boundary}
\label{sec:validation-proof}
The mathematical proof kernel is the function \texttt{c}. It is copied
without a syntax-tree change from the 897-byte control. The simplification
is entirely before that kernel: establish the exact raw data language once,
then specialize the public-goal validator to that established domain.

\subsection{Recognizing integer/list trees without a recursive type mode}
Let $\mathcal J$ be the finite values returned by the standard Python JSON
decoder. Its possible node types are exact integers, exact lists, strings,
dictionaries, Booleans, floating-point numbers, and \texttt{None}.
The parser's additional \texttt{NaN} and infinity constants, when accepted,
are floating-point numbers. Let $\mathcal L\subseteq\mathcal J$ be the
least class containing every exact integer and every finite list whose
elements lie in $\mathcal L$.

Let $\Sigma$ consist of the ten decimal digits, left and right square
brackets, comma, space, and minus sign. Define
\[
 R(x)\quad\Longleftrightarrow\quad
 \texttt{repr}(x)\text{ contains only characters in }\Sigma.
\]
In Python, \texttt{not repr(x).strip('0123456789[], -')} tests precisely
this condition: removing allowed characters from both ends leaves an
empty string exactly when no forbidden character occurs anywhere.

\begin{lemma}[Exact raw-domain recognition]
\label{lem:repr}
For every finite JSON-decoded value $x$ for which the relevant operations
finish, $R(x)$ holds if and only if $x\in\mathcal L$.
\end{lemma}
\begin{proof}
An exact integer has a decimal representation with an optional minus sign.
A list representation adds only square brackets, commas, and spaces to
the representations of its elements. Structural induction therefore proves
the forward construction $x\in\mathcal L\Rightarrow R(x)$.

For the converse, every disallowed node contributes a forbidden character
which an enclosing list preserves. A string contributes a quote character;
a dictionary contributes a brace; \texttt{True}, \texttt{False}, and
\texttt{None} contribute letters. A finite floating-point representation
contains a decimal point or an exponent marker; infinities and NaNs contain
letters. Consequently no tree containing any such node satisfies $R$.
The argument concerns values produced by the standard JSON decoder, not
arbitrary Python objects with user-defined representation methods.
\end{proof}

The new guard applies to the \emph{whole decoded input before popping the
goal}. Thus an unused prefix cannot hide an invalid value. Top-level
integers may pass the raw-domain test, but fail the subsequent list operation;
an empty top-level list similarly fails when its goal is popped. Exceptions
reject, as in the inherited protocol.

\subsection{Specializing the scoped goal validator}
The old validator \texttt{V(a,d)} has two modes: $d=0$ recursively checks
raw integer/list types, while $d\ge2$ checks public syntax at opener depth
$d-2$. The new \texttt{v(a,d)} has only the latter task and starts directly
at opener depth $0$.

On the already established domain $\mathcal L$, multiplying a node by zero
distinguishes the cases: integers yield integer zero and lists yield an
empty list. Thus \texttt{a * 0 == 0} selects exactly the numerical-leaf case.
For an integer leaf, the new scoped condition is
\[
 d\bmod2=1\quad\text{and}\quad
 \bigl(a=0\ \text{or}\ (2<a<d\text{ and }a\bmod2=1)\bigr).
\]
For a list, the permitted arities are $1$ and $2+d\bmod2$, together with a
diagonal triple in the full native profile. Its children are checked at
depth $d+5-\lvert a\rvert$. This is the inherited arity convention, including
its dummy protected slots; it is not a new binding discipline.

\begin{lemma}[Depth-shift specialization]
\label{lem:depthshift}
For each $a\in\mathcal L$ and nonnegative opener depth $d$, the new
\texttt{v(a,d)} successfully recognizes the same scoped public syntax as
the old \texttt{V(a,d+2)}. The Eq/D restriction satisfies the analogous claim
after removing the public diagonal-Q alternative from both validators.
\end{lemma}
\begin{proof}
Induct on the input tree. Adding two to depth preserves parity; the old
bound $a<(d+2)-2$ becomes $a<d$. Raw-mode alternatives are false because
the old depth is positive. At a list, both validators impose the same arity
condition. The old child depth is $(d+2)+5-\lvert a\rvert$, exactly two
greater than the new child depth, so the induction hypothesis applies.
For an empty list, a short-circuited diagonal access may raise rather than
return false. This is rejection in both checker contracts. The statement
is about successful recognition, not total Boolean functions on all objects.
\end{proof}

\begin{theorem}[Preservation of the validated presentation]
\label{thm:preservation}
Under the finite-certificate mathematical interpretation of the programs,
the revised full native checker and its 897-byte control recognize the
same valid public proof certificates. The corresponding Eq/D programs
also agree. Hence each revised checker presents the same native consequence
relation and inherits the same whole-domain bi-interpretation with iPRA.
\end{theorem}
\begin{proof}
The old raw validator recognizes exactly $\mathcal L$. The old split check
on the remaining prefix and canonical goal is equivalent to requiring that
the original whole input belong to $\mathcal L$ and that the goal have the
required canonical syntax. Lemma~\ref{lem:repr} establishes the former
condition in the new program; Lemma~\ref{lem:depthshift} establishes the
latter. Both then pop the same goal, expose the same packet stream in the
same order, and execute the unchanged kernel with the same initial context
and freshness bound. All successful kernel steps therefore coincide.
No axiom, introduction mode, scope check, or inference capability is added.
The inherited soundness and completeness theorem transfers to the revised
program. Composing the inherited native/ordinary and ordinary/iPRA
interpretations gives the final claim.
\end{proof}

\paragraph{Operational limits and diagnostics.}
The theorem does not assert identical behavior at every Python resource
limit. The old raw traversal uses Python recursion; the new representation
test has different implementation limits. An existing valid computation
certificate with an unused integer/list prefix nested 600 levels is accepted
by the new program but exhausts the old raw validator's recursion limit in
the tested environment. Both still enforce the same mathematical data
language. Conversely any resource failure rejects, and no finite-hardware
program accepts every arbitrarily large finite proof. The error-message
text is retained verbatim, but malformed values can reach a different
exception class. The acceptance convention remains normal exit with exactly
\texttt{ACCEPT} followed by a newline.

\clearpage
\section{The complete standalone implementation}
\label{sec:implementation}
The source below is the entire trusted checker apart from the Python runtime
and its standard JSON module. The selected program imports no producer,
axiom database, reference implementation, or test module.
\lstinputlisting{code/native_typed_full.py}

The four packet fields are $[m,s,u,t]$, read from right to left in the outer
input list. The checker first verifies the major $m$. A matching equation
major supplies transport; a matching diagonal Q major supplies QE. A
proved D major exports its closure under a fresh parameter. A neutral major
with truthy $u$ supplies a computation equation; a neutral major with falsy
$u$ supplies an introduction. The falsy values $[]$ and $0$ distinguish the
case and equation-progress choices. Only a literal equation obtains the
progress assumption. The exact packet conventions and the inherited
certificate decoder are included with the sources.

Certified zero instances and the syntactic exclusion of a literal-zero
transport source are inherited implementation devices. Their admissibility
is not a consequence of the new representation lemma. The required
mathematical arguments are retained in
Appendices~\ref{sec:arity-checker} and~\ref{zs:section}; the imported
source snapshots identify the current matched control. New certificates
need no translation to use the revised front end.

\subsection{Why a tempting additional reduction was rejected}
At entry to the local checker, one has $k=1-\lvert g\rvert$, suggesting
that $k$ could be removed as a separately transmitted argument. That alone
does not eliminate its role. The nested opener closes over the local $k$,
which is decremented \emph{before} the final goal and hypothesis openings.
Those calls must admit the newly allocated name.

Replacing that changing bound by one fixed expression gives a smaller
program which accepts simple computation certificates but rejects existing
contextual-induction and transport-under-binder certificates. A correct
version which computes the entry value from $g$ and retains the decrement
is eight bytes larger than the original. The regression receipt is included
as evidence of a rejected attempt; no incorrect experimental checker is
part of the delivered code. The context-length invariant itself was already
recorded in the earlier local-scopes package.

\clearpage
\section{Measurements and validation}
\label{sec:measurements}
\subsection{Measurement convention}
Measurements use CPython 3.12.14 and zlib 1.3.2. Source is parsed, docstrings
are removed, and \texttt{ast.unparse} is used with one final newline. The
principal coordinates are UTF-8 bytes, significant Python tokens, and AST
nodes. Significant tokens omit encoding, layout, comments, and the end
marker. Gzip uses compression level 9 and a fixed zero timestamp. All source
files measured here include their input and output interfaces. The original
diagnostic sentence has not been shortened to manufacture a byte saving.

\begin{center}
\begin{tabular}{lrrrr}
\toprule
Complete standalone program & Bytes & Tokens & AST & Gzip-9\\
\midrule
Full Eq/Q/D control & 897 & 401 & 468 & 505\\
Full Eq/Q/D refinement & 870 & 383 & 442 & 507\\
Eq/D control & 878 & 390 & 454 & 498\\
Eq/D refinement & 849 & 370 & 426 & 499\\
\bottomrule
\end{tabular}
\end{center}


For full Eq/Q/D, the reductions are 3.01\% in bytes, 4.49\% in tokens,
and 5.56\% in AST nodes. For Eq/D they are 3.30\%, 5.13\%, and 6.17\%.
The compressed programs grow by two bytes and one byte, respectively;
compression was not the optimization objective. The Eq/D profile changes
only the public formula grammar: Q remains an intermediate resource and is
definable publicly by D with an unused witness variable.

The optional lexical variant filters the raw input text before JSON parsing,
using the digits, brackets, comma, minus sign, and the four JSON whitespace
characters. The JSON parser checks the placement of these characters.
It has the same intended type boundary and a direct lexical proof, but at
885 normalized bytes it is larger than the selected 870-byte representation
variant. It is included for comparison rather than called a second record.

\subsection{What the executable checks establish}
The reproduced comparison has zero discrepancies. For each public profile, it checks \textbf{21,024} JSON-derived values against an independent integer/list predicate and performs \textbf{240,120} old/new scoped-validator comparisons. The 101 input fixtures give 21 acceptances and 80 rejections for full Eq/Q/D, and 16 acceptances and 85 rejections for Eq/D, matching their controls. Eleven disallowed unused-prefix types are rejected in ordinary and optimized Python execution. The corpus includes one metadata-shaped JSON object as a deliberate top-level rejection. The proof-kernel AST is unchanged in both profiles.


The proof corpus exercises computations, contextual equation induction,
transport beneath D and Q, closure use, compound case introduction,
zero-source recovery, and rejection of unproved majors and ill-scoped
names. Additional adversarial cases put disallowed JSON values into unused
prefixes, where a goal-only check would miss them. The validator comparisons
include arrays of incorrect arities, empty nodes, aliases, and bound-slot
limits. Runs with \texttt{python3 -O} verify that acceptance does not rely
on assertions.

Finite tests support the implementation theorem; they do not prove it or
constitute a Lean verification. The original incoming package's supplied
reproduction was also run successfully before changing anything: 642 process
invocations, 18,060 scoped-validator comparisons, and 3,834 raw-type
comparisons, all with the expected outcomes. Those are baseline receipts,
separate from the new checks reported above. The current Lean sources were
inspected but not rebuilt in this work, and no Lean/Python correspondence
theorem is claimed.

\Needspace{240pt}
\section{Further research questions}
\label{sec:further}
\begin{enumerate}
\item Can the mathematical core be simplified while also reducing a fully
validated checker under a fixed certificate contract? A proposed removal
must recover the deleted rule and preserve both interpretation composites;
shorter surface syntax by itself does not settle the question.
\item Can the arity decoder, public validator, and successful-trace invariant
be formalized together, then connected to the existing Lean derivability
relation? This would address the main remaining trust gap for the Python
implementation without changing iPRA's mathematical presentation.
\item Can numerical terms serve as atomic zero assertions, replacing an
equation $s=t$ by a fixed scalar distance term $\delta(s,t)=0$, at lower
\emph{total} trusted cost? The fixed expansion of $\delta$ and the internal
reflection and round-trip proofs must be counted; adding an uncharged
equality-test primitive would merely move complexity.
\item What restricted class of representations and implementations would
make a meaningful lower-bound theorem possible? The present collection
gives best observed coordinates in particular protocols, not a
representation-independent lower bound for axiomatizations.
\end{enumerate}

\section*{Artifact guide and attribution}
\addcontentsline{toc}{section}{Artifact guide and attribution}
The archive contains the compiled report, its complete expanded LaTeX source,
the editable component sources, full and Eq/D standalone checkers, an annotated
implementation, a lexical alternative, matched frozen controls,
proof examples and adversarial cases, reproducible measurements and checks,
and source provenance. Run \texttt{python3 reproduce.py} for the executable
evidence and \texttt{python3 build\_report.py} to compile the report with
\texttt{pdflatex}.

The mathematical construction, certificate kernel, and inherited appendix
proofs come from the pinned ProveIt development under its MIT No Attribution
license. Their attribution is retained nonetheless. The contribution of
this note is the smaller validation boundary and its correctness comparison.
The following appendices retain their source terminology, including words
such as ``new'' relative to their original reports; those words do not
reclassify the inherited results as new work here. Parallel reviews in the
evidence folder are assistant reviews, not external human reviews.

\clearpage
\appendix
% New mathematical result: complementary introduction choices.
% Requires amsmath/amssymb/amsthm and theorem, lemma, proposition environments.
% The inherited arithmetic-to-iPRA bridge must be supplied separately.

\section{Complementary introductions with one local assumption}
\label{sec:choice-introduction}

The old paired introduction retains both the zero instance and an
induction instance when its matrix is an equation. Both together are
unnecessary. A certificate may choose either one, provided the induction
choice is restricted to a literal equation. This changes neither the
formula language nor the meaning of a formula. The omitted combination
is recovered by a short nested introduction--elimination construction.

\subsection{Complete native syntax and rule inventory}

The scalar numerical language is
\[
 t::=x\mid0\mid St\mid I(n,a;u.h).
\]
Only $h$ binds the state variable $u$; $n$ and $a$ do not.
The update may contain fixed numerical parameters. The computation
axioms in every retained context are
\[
 I(0,a;u.h)=a,\qquad
 I(Sn,a;u.h)=h[I(n,a;u.h)/u].
 \tag{I0, IS}\label{eq:choice-computations}
\]
The argument also applies unchanged to the full fixed-description PR
language with its projection, composition, and recursion equations.
No numerical variable ranges over function descriptions, and no uniform
evaluator is postulated.

Native formulas and reusable context clauses are
\[
 A::=(s=t)\mid\mathsf D_{x,y}A,\qquad
 \langle A;y.B(y)\rangle,\quad y\notin\operatorname{FV}(A).
\]
Bound variables are renamed before capture-free substitution.
The two variables $x,y$ in a paired formula are distinct.
An ordinary assumption $F$ is the constant clause
$\langle H;v.F\rangle$, where $v$ is absent from $F$ and $H$ is the
fixed computation theorem $I(0,0;u.0)=0$. Thus ordinary hypothesis
use is a clause Use preceded by this computation proof.

There are computation axioms and exactly the following inference
patterns:
\[
\frac{\Gamma\vdash A}{\Gamma\vdash B(t)}
\quad\mathsf{Use}
\quad(\langle A;y.B(y)\rangle\in\Gamma),
\]
\[
\frac{\Gamma\vdash\mathsf D_{x,y}B(x,y)\qquad
      \Gamma,\langle B(a,0);y.B(a,y)\rangle\vdash F}
     {\Gamma\vdash F}
\quad\mathsf{DE},
\qquad
\frac{\Gamma\vdash s=t\qquad\Gamma\vdash F(s)}
     {\Gamma\vdash F(t)}
\quad\mathsf{Tr},
\]
\[
\frac{\Gamma,B(t,\eta_\varepsilon(z))\vdash B(t,Sz)}
     {\Gamma\vdash\mathsf D_{x,y}B(x,y)}
\quad\mathsf{DI}_{\varepsilon},
\qquad
\eta_0(z)=0,\quad\eta_1(z)=z.
\tag{Choice DI}\label{eq:choice-di}
\]
Here $\varepsilon$ is a certificate choice in $\{0,1\}$, and
$\varepsilon=1$ is permitted only when $B$ has a literal equation
root. An equation beneath a binder is not a literal equation.
Call $\varepsilon=0$ the \emph{case choice} and $\varepsilon=1$
the \emph{progress choice}. Both retain exactly one ordinary
formula assumption. Compound matrices have only the case choice.

In DI the eigenparameter $z$ is fresh for $\Gamma$, the witness
$t$, and every fixed parameter of $B$. In DE the eigenparameter
$a$ is fresh for $\Gamma$, the conclusion $F$, and every fixed
parameter of $B$. In transport, $F(v)$ is any fixed native
formula template. Displayed local assumptions and clauses are
discharged. There are no primitive reflexivity, coverage,
separation, ordinary logical, or additional induction rules.

We call this calculus $\mathcal C$. The original comparison
calculus $\mathcal N$ has identical syntax, computation,
Use, DE, and transport, but its introduction is
\[
\frac{\Gamma,B(t,0),[B(t,z)]_{\mathrm{at}}\vdash B(t,Sz)}
     {\Gamma\vdash\mathsf D_{x,y}B(x,y)}.
\tag{Old DI}\label{eq:choice-old-di}
\]
The bracket supplies $B(t,z)$ exactly at a literal equation.

\subsection{Exact equivalence and the unused-witness construction}

\begin{lemma}[Structural admissibility]
\label{lem:choice-structural}
Weakening, exchange, contraction, renaming, ordinary-assumption cut,
and capture-free numerical substitution are admissible in $\mathcal C$.
\end{lemma}
\begin{proof}
Induct on derivations. Rename local eigenparameters before inserting
proofs or terms. Numerical substitution preserves a formula's root
constructor, so an allowed progress choice remains allowed.
The selected choice does not change under substitution or cut.
For ordinary-assumption cut, replace each use of its constant clause
by the supplied proof. Computation schemas are closed under
capture-free numerical substitution. These facts verify every case.
\end{proof}

\begin{theorem}[Identity-language equivalence]
\label{thm:choice-equivalence}
For every retained clause context $\Gamma$ and native formula $F$,
\[
 \Gamma\vdash_{\mathcal C}F
 \quad\Longleftrightarrow\quad
 \Gamma\vdash_{\mathcal N}F.
\]
The transformations are effective on finite derivations and leave
all numerical and formula syntax unchanged.
\end{theorem}
\begin{proof}
Each new introduction is an instance of the old one after weakening
by the omitted local assumption. For the case choice at an equation,
add the induction instance. For the progress choice, add the zero
instance. Compound case introduction is already the old rule.
Thus every new derivation translates to an old one.

For the reverse direction, transform an old derivation recursively.
Only its atomic DI case needs a construction. Suppose its matrix is
the equation $E(x,y)$ and its transformed premise is
\[
 \pi:\quad\Gamma,E(t,0),E(t,z)\vdash E(t,Sz).
\]
The outer new inference will be case DI with witness $t$. It is
therefore enough to prove $E(t,Sz)$ under
$\Gamma,E(t,0)$.

Rename the eigenvariable in $\pi$ to a fresh $v$. Choose a new
variable $u$ absent from $E(t,y)$ and all retained data. Under
$\Gamma,E(t,0)$, progress DI with witness zero introduces
\[
             \mathsf D_{u,y}E(t,y)
\]
from exactly the renamed premise $\pi[v/z]$.
The bound coordinate $u$ is deliberately unused.

Apply DE to this paired formula. Because $u$ is unused, the opened
clause is exactly
\[
               \langle E(t,0);y.E(t,y)\rangle;
\]
it contains no new existential witness. The retained ordinary
assumption $E(t,0)$ supplies its antecedent. Use at $Sz$ gives
$E(t,Sz)$, and DE closes. Finally the outer case DI introduces
the original $\mathsf D_{x,y}E(x,y)$ with witness $t$.
The original premise proof $\pi$ has been inserted once.

An old compound introduction already is the new case choice.
All other rule cases are unchanged. This completes the induction
on the original derivation.
\end{proof}

The unused inner coordinate is essential. Introducing
$\mathsf D_{x,y}E(x,y)$ internally and then eliminating it would
give an arbitrary fresh witness, not necessarily the original
$t$. The construction instead keeps $t$ as a fixed parameter
of $E(t,y)$ and binds only the unused coordinate $u$.

\begin{proposition}[Linear bound on inference nodes]
\label{prop:choice-node-bound}
Count each computation axiom, Use, DI, DE, and transport as one
inference node, expanding ordinary-assumption use as a computation
of $H$ followed by Use. If an old derivation has $N$ nodes and
$r$ atomic DI nodes, the displayed transformation has exactly
\[
                   N+5r\leq 6N
\]
nodes.
\end{proposition}
\begin{proof}
Each atomic old DI node is replaced by six nodes: outer case DI,
inner progress DI, DE, Use of the opened clause, Use of the
retained base assumption, and the computation of $H$ needed for
that last ordinary-assumption use. Its old premise subtree
occurs once. All other inference nodes are copied once.
\end{proof}

This counts inference-tree nodes. It is not a bound on serialized
JSON bytes, formula sizes, checker runtime, or the number of
consumed packet cells when hypothesis closure is implicit.
In the formula-context Q/D presentation, the corresponding
replacement has outer case DI, inner progress DI, DE, QE,
and two Hyp nodes, giving the same six-node count.

\subsection{Induction and numerical cases separate cleanly}

\begin{lemma}[Contextual equation induction]
\label{lem:choice-ei}
The progress choice, DE, and Use derive
\[
\frac{\Gamma\vdash E(0)\qquad
      \Gamma,E(z)\vdash E(Sz)}
     {\Gamma\vdash E(t)}
\tag{EI}\label{eq:choice-ei}
\]
for every equation $E$, with the usual freshness condition.
\end{lemma}
\begin{proof}
Choose an unused witness coordinate $u$. Progress DI introduces
$\mathsf D_{u,y}E(y)$ from the given step. DE opens the clause
$\langle E(0);y.E(y)\rangle$. Supply its antecedent using the given
base proof and use the clause at $t$. Close DE.
\end{proof}

\begin{lemma}[Cases, equality, and binder congruence]
\label{lem:choice-preinduction}
The case choice derives, for arbitrary native formulas,
\[
\frac{\Gamma\vdash F(0)\qquad\Gamma\vdash F(Sz)}
     {\Gamma\vdash F(t)}
\tag{NC}\label{eq:choice-nc}
\]
with the usual freshness condition. Equality reflexivity,
symmetry, transitivity, numerical congruence, and uniform
matrix interderivability under $\mathsf D$ are derivable
without using the progress choice.
\end{lemma}
\begin{proof}
For cases, use case DI with an unused witness coordinate and
matrix $F(y)$, ignoring its retained base in the successor
premise. DE opens $\langle F(0);y.F(y)\rangle$; the given zero
proof and Use at $t$ finish.

For reflexivity put $r=I(0,t;u.0)$. Use the computation equation
$r=t$ as both transport premises with motive $v=t$ to obtain
$t=t$. The usual explicit equality motives give symmetry,
transitivity, and numerical congruence.

For binder congruence, open the old paired formula at a fresh
witness $a$ and introduce the new one by case DI at $a$.
Convert its retained base to the old base, use the opened
old clause at the successor, and convert back. Ordinary cut
closes the construction, and the reverse conversion gives
the converse. No induction is used.
\end{proof}

The primitive progress choice is still restricted to equations.
Allowing it at arbitrary compound matrices would, by the
argument for \eqref{eq:choice-ei}, supply induction for
arbitrary formulas and change the intended theory.

\subsection{The entire ordinary logical interface}
\label{subsec:choice-logic}

Here is a predecessor-free reconstruction of the ordinary logical
operations. It makes explicit that the case choice supplies the
logical interface independently of progress induction.

Let $C(v,a,b):=I(v,a;u.b)$ with $u$ unused. Computation gives
$C(0,a,b)=a$ and $C(Sv,a,b)=b$. Define the gate on fixed finite
formula syntax by
\[
 G_v(s=t):=(C(v,0,s)=C(v,0,t)),\qquad
 G_v(\mathsf D_{x,y}B):=\mathsf D_{x,y}G_vB,
\]
renaming bound variables away from the control. Equality,
binder congruence, and external structural induction give
\[
 \vdash G_0A,\qquad G_{Sv}A\dashv\vdash A,\qquad
 A\dashv\vdash\mathsf D_{u,y}G_yA
 \quad(u,y\notin\operatorname{FV}(A)).
\tag{Gate and padding}\label{eq:choice-gate}
\]
For padding, case DI proves the forward direction. DE, the
zero-gate theorem, and Use at $1$ prove the reverse direction.

Define a selector between two arbitrary formulas by
\[
 M_v(s=t,u=w):=(C(v,s,u)=C(v,t,w))
\]
and, with aligned fresh paired variables, by
\[
\begin{aligned}
 M_v(\mathsf D_{x,y}A,\mathsf D_{x,y}B)
   &:=\mathsf D_{x,y}M_v(A,B),\\
 M_v(E,\mathsf D_{x,y}B)
   &:=\mathsf D_{x,y}M_v(G_yE,B),\\
 M_v(\mathsf D_{x,y}A,E)
   &:=\mathsf D_{x,y}M_v(A,G_yE),
\end{aligned}
\]
where $E$ is an equation. Every call decreases maximum input
binder depth. Padding and congruence prove
$M_0(A,B)\dashv\vdash A$ and
$M_{Sv}(A,B)\dashv\vdash B$.

Put $\delta(v)=C(v,1,0)$ and $p(z)=I(z,0;u.\delta(u))$.
Only the fixed computations $p(1)=1$ and $p(2)=0$ are needed.
The full logical macro inventory is
\[
\begin{aligned}
 A\Rightarrow_*B&:=\mathsf D_{u,z}M_z(A,B),\\
 A\wedge_*B&:=\mathsf D_{u,z}G_zM_{p(z)}(A,B),\\
 \exists_*x\,A(x)&:=\mathsf D_{x,y}G_yA(x),\\
 A\vee_*B&:=\exists_*x\,M_x(A,B),\\
 \forall_*x\,A(x)&:=
        A(0)\wedge_*\mathsf D_{u,z}G_zA(z),\\
 \bot_*&:=(0=1).
\end{aligned}
\tag{Logical macros}\label{eq:choice-macros}
\]
Newly bound variables are fresh; $u$ is unused.

Implication introduction obtains its assumption from the zero
selector and supplies the consequent at successors, then uses
case DI. Elimination opens its clause and uses it at $1$.
Existential introduction uses its witness and the successor
gate; elimination opens at a fresh witness, supplies the
zero-gate theorem, and uses the clause at $1$.

Disjunction introduction uses witnesses $0$ and $1$.
For elimination, discharge the two branch proofs as
$A\Rightarrow_*F$ and $B\Rightarrow_*F$. NC applied to
$M_z(A,B)\Rightarrow_*F$ gives this implication at every
control: the two endpoint cases follow from the branch
implications. Existential elimination provides $M_a(A,B)$
and the implication gives $F$.

Given $A$ and $B$, selector endpoints and NC give $M_v(A,B)$
at an arbitrary control. Substitute $v=p(Sz)$ and use the
successor gate for conjunction introduction. For projections,
open its zero-gated clause and specialize at $1$ and $2$;
the two fixed computations give $B$ and $A$, respectively.

These rules prove existential numerical coverage:
\[
     (x=0)\vee_*\exists_*v\,(x=Sv).
\tag{Coverage}\label{eq:choice-coverage}
\]
Use NC. Zero takes the left injection of reflexivity, and a
successor takes the right injection with the evident witness.
Coverage gives equality cases even when the retained context
contains $x$; its elimination adds $x=0$ or $x=Sv$ rather
than silently changing that retained context.

Universal introduction follows from a fresh-parameter proof
by substituting zero and successors and forming the two
conjuncts. For elimination, take its $A(0)$ conjunct and
open the paired conjunct. Its base is a zero-gate theorem.
Coverage of the target chooses the zero instance from
$A(0)$ or a successor instance from the opened clause
and gate removal.

Finally, under $0=1$, computation and transport give
$s=C(0,s,t)=C(1,s,t)=t$ for arbitrary numerical terms.
Every equation follows, and external induction on formula
syntax with case DI gives every native formula. This
supplies falsity elimination. No excluded middle for
quantified formulas has entered the construction.

\subsection{Whole-domain realizations and both formula composites}

For a standard first-order reading, let $T_{\rm eq}$ have
intuitionistic equality logic, the same numerical signature
and computation equations, numerical coverage, separation
$(0=1)\to\bot$, and contextual equation induction.
Its interpretation of the native syntax is
\[
 (s=t)^\circ=(s=t),\qquad
 (\mathsf D_{x,y}B)^\circ
   =\exists x\bigl(B^\circ(x,0)\to\forall y B^\circ(x,y)\bigr),
\]
\[
 \langle A;y.B\rangle^\circ
       =A^\circ\to\forall y B^\circ.
\]
The domain is the entire numerical domain and interpreted
equality is actual equality.

Case DI is sound in this realization by numerical coverage.
Progress DI is sound by contextual equation induction under
the conditional base. Use, DE, transport, and computation
have their displayed ordinary meanings. Conversely, the
native macros derive every ordinary logical rule,
\eqref{eq:choice-ei} derives equation induction, and
\eqref{eq:choice-coverage} derives coverage. Separation
translates to the identity implication
$(0=1)\Rightarrow_*(0=1)$.

For the essential paired formula composite,
\[
 \mathsf D_{x,y}B
 \dashv\vdash
 \exists_*x\bigl(B(x,0)\Rightarrow_*\forall_*y\,B(x,y)\bigr).
\tag{Paired composite}\label{eq:choice-paired-composite}
\]
Forwards, DE opens its clause at a fresh $a$; under $B(a,0)$
that clause supplies $B(a,y)$ at a fresh arbitrary $y$.
Introduce the universal, implication, and existential
macros and close DE. Backwards, eliminate the existential
at a fresh $a$. Introduce the paired formula with witness
$a$ by case DI: the retained base, supplied implication,
and universal elimination at $Sz$ give its successor premise.

The ordinary interpretations of the six logical macros
have their ordinary round trips. For conjunction, its true
base and the instances $1,2$ recover both operands;
the converse uses selector cases. The universal has all
successor instances in its paired conjunct and zero in
its first conjunct; coverage gives exactly the ordinary
universal. Implication, disjunction, and existential
quantification use their selector or gate endpoints.
Falsity uses the explicit ordinary separation bridge.
External induction on finite formulas, together with
\eqref{eq:choice-paired-composite}, gives both formula
composites. Clauses correspond in both directions to
$A\Rightarrow_*\forall_*yB$ by the derived logical rules.

The numerical composites are literal identities with
internally provable graph $x=y$. No standard-numeral
restriction, retract, quotient, or equality reflection
assumption is being inserted at this stage. The separately
supplied arithmetic bridge identifies $T_{\rm eq}$ with
full iPRA; the present change preserves that bridge
unchanged because Theorem~\ref{thm:choice-equivalence}
preserves every native sequent.

\subsection{What the theorem does and does not claim}

The new inference interface retains one ordinary assumption
at every paired introduction. It preserves the entire
consequence relation, not just equations or closed
theorems. Consequently the established full-iPRA
interpretations and equiconsistency transfer unchanged.

This result is different from retaining only progress
introduction and strengthening closure elimination.
Here the original elimination rule is retained, and
the case and progress introductions coexist. It is also
different from allowing arbitrary compound progress,
which would supply stronger induction.

A complete-checker comparison must still count the
certificate choice, its equation-root restriction,
substitution, context handling, computation, and the
public input/output interface. The mathematical theorem
alone is not a byte-count or minimality theorem.
The inference-node bound is distinct from source-size
measurements and from serialized certificate size.
The proof is a written universal proof; finite regression
tests do not constitute a Lean/Rocq formalization of it.

\section{Arithmetic strength and both ordinary interpretation composites}
\label{sec:ordinary-bridge}

We spell out the arithmetic bridge to avoid identifying standard numerical
agreement with bi-interpretability. First elaborate every fixed iterator
into an ordinary function symbol. At stage zero the symbols are $0,S$.
Given a fixed earlier-stage term $h(u,\vec p)$, introduce
$i_h(n,a,\vec p)$ with the two equations
\[
i_h(0,a,\vec p)=a,\qquad
i_h(Sn,a,\vec p)=h(i_h(n,a,\vec p),\vec p).
\]
Only finite descriptions occur; the domain contains no code for a
universal evaluator. Formal parameters are distinct when a description
is declared, while actual arguments can coincide.

Let $T_I$ be ordinary intuitionistic first-order logic with actual
equality and genuine falsity, these staged equations, contextual
induction for literal equations, and
\[
\forall x\,(x=0\vee\exists y\,x=Sy),
\qquad (0=S0)\to\bot.
\]
The second sentence is the falsity bridge. It must be stated: positive
numerical equations by themselves admit a singleton structure. Native
explosion from $0=S0$ and a translation of ordinary falsity to that
equation do not authorize silently dropping the bridge from $T_I$.

The native-to-ordinary translation is
\[
(s=t)^\circ=(s=t),\qquad
(\mathsf D_{x,y}B)^\circ=
\exists x\bigl(B^\circ(x,0)\to\forall y B^\circ(x,y)\bigr).
\]
A clause $\langle A;y.B\rangle$ becomes $A^\circ\to\forall yB^\circ$.
For the auxiliary closure used by the compact checkers,
$(\mathsf Q_yB)^\circ=B^\circ(0)\to\forall yB^\circ$.
The interpretation domain is all elements, defined by $x=x$;
equality is actual equality.

Case DI is sound because numerical coverage combines its retained base
with its successor proof. Progress DI is sound by literal-equation
induction under a temporarily assumed base. Use, DE and transport are
ordinary logical inferences under the translations. Conversely the
derived logical operations, equation induction, cases and explosion in
Section~\ref{sec:choice-introduction} recover the rules and axioms of
$T_I$. An auxiliary closure is interderivable with
$\mathsf D_{v,y}B(y)$ for unused $v$: in one direction use case DI with
the closure and its base; in the other eliminate the paired formula
and return its supplied closure. Thus formulas and context clauses have
the stated meanings in both presentations.

\subsection{The arithmetic dependency order}

The scalar update has no recursion-counter argument. We do not conceal
one in an encoding. The inherited internal proofs, reproduced in the
appendices, establish the required arithmetic in this order:
\begin{enumerate}
\item Using equation induction and numerical cases, construct a fixed
scalar predecessor term and prove $P0=0$ and $P(Sn)=n$.
Existential coverage then gives $x=0\vee x=S(Px)$: in the successor
case $x=Sy$, apply $P$ to obtain $Px=y$ and substitute.
\item Construct truncated subtraction and a numerical characteristic
of actual equality. Prove its reflection on the entire original domain.
\item Compile each fixed quantifier-free formula into its characteristic
equation and recover every contextual quantifier-free induction instance.
\item With that induction available, construct pairing and projections
by a bounded square-root search and compile every fixed general
primitive-recursive description into a scalar iterator term.
\end{enumerate}
No pairing or faithful equality detector is assumed to obtain
predecessor. No appeal to standard-model correctness replaces an
internal equality. Every induction fixes its passive parameters;
diagonal substitution occurs only into a previously proved identity.

For a target recursion
\[
f(0,\vec p)=g(\vec p),\qquad
f(Sn,\vec p)=h(n,f(n,\vec p),\vec p),
\]
let $\pi,L,H$ be the proved pairing and projections. The compiler uses
\[
V(n,\vec p)=I\bigl(n,\pi(0,g^*(\vec p));
u.\pi(SL(u),h^*(L(u),H(u),\vec p))\bigr),
\qquad f^*(n,\vec p)=H(V(n,\vec p)).
\]
Equation induction proves $L(V(n,\vec p))=n$; the two projection
identities then prove the target defining equations. The construction
is external recursion on each finite PR description, not an internal
quantification over arbitrary functions.

\begin{theorem}[Whole-domain bi-interpretability and equiconsistency]
The ordinary realization $T_I$ and full first-order iPRA are
bi-interpretable. In both theories $x=y$ defines the isomorphism from
the original domain to its double interpretation. The native Choice-DI
calculus is therefore a faithful proof presentation of a theory
bi-interpretable and equiconsistent with iPRA.
\end{theorem}
\begin{proof}
Interpret an iterator in iPRA by primitive recursion that ignores its
counter argument in the update. iPRA proves its equations, literal
equation induction, coverage and separation. In the reverse direction
use the scalar compiler. The predecessor and equality-reflection proofs
supply the successor axioms and every quantifier-free induction instance.
All ordinary logical rules are preserved.

For two iterators with the same seed and provably equal updates,
equation induction on the equality of their values proves uniqueness at
every counter, with all parameters fixed. The corresponding argument
applies to ordinary primitive recursion. Induct externally on a finite
description: its double translation satisfies its original defining
equations, so internal uniqueness identifies the two functions. This
also proves coherence when substitution changes the staged syntactic
description of an update.

The domain and equality are unchanged; hence the identity graph is
internally total, bijective and preserves every interpreted operation.
The displayed ordinary meanings of the logical macros and the paired
binder give both formula composites by external induction on finite
formula syntax. Clauses are similarly equivalent to their conditional
universal meanings. The explicit falsity bridge preserves contradiction
in both directions. Effective translations of finite proofs give
equiconsistency, without a claim that either theory proves its own
consistency.
\end{proof}

% Adapted with attribution from the pinned unparenthesized one-hole archive.
\section{Internal scalar predecessor bootstrap}
\label{sec:hybrid-bootstrap}

This scalar bootstrap is inherited from the supplied scalar development~\cite{erasure} and is reproduced to keep the scalar checker self-contained. It is not a new contribution of this paper. The proof uses a factor map fixing odd states and a saturated-iteration construction of predecessor.

The available arithmetic terms are $0$, $S$, and
\[
 I(n,a;u.h),\qquad
 I(0,a;u.h)=a,\qquad
 I(Sn,a;u.h)=h[I(n,a;u.h)/u].
\]
The update binds only $u$.  All other variables in it are fixed numerical
parameters.  We use equality reasoning, the already reconstructed
zero/successor cases, and induction on a single literal equation.
Every induction below keeps its other parameters fixed.  Predecessor,
successor injectivity, pairing, and induction on arbitrary
quantifier-free formulas are not assumptions of the argument.

\subsection{Elementary terms and two finite cycles}

Write $1=S0$ and use larger numerals as successor abbreviations.  Define
\[
 C(n,a,b)=I(n,a;u.b),\qquad
 Z(n)=C(n,1,0),\qquad
 a+b=I(b,a;u.Su).
\]
The update variable of $C$ is absent from its body.  Consequently
\[
 C(0,a,b)=a,\qquad C(Sn,a,b)=b.
\]
Numerical cases show that $Z(n)$ is $0$ or $1$.
The usual commutative-monoid equations for addition follow by literal
equation induction.  For example, $0+b=b$ and $Sa+b=S(a+b)$ are
proved by induction on $b$; these yield commutativity by a further
equation induction, and associativity follows by induction on the last
summand.  We use these equations to reorder finite sums.  The notation
$2a$ means $a+a$ and $3a$ means $a+a+a$; no multiplication operation
is needed for this bootstrap.

Put
\[
 N(n)=I(n,0;u.Z(u)),\qquad
 j(u)=Z(u)+2N(u),\qquad
 M(n)=I(n,0;u.j(u)).
\]
The range of $N$ is contained in $\{0,1\}$: its zero value is $0$,
and its successor value is a value of $Z$.  This uses numerical cases
on its counter, not induction on a disjunction.  Hence
\begin{equation}
 N(n)+Z(N(n))=1,\qquad
 Z(Z(N(n)))=N(n),\qquad N(n+2)=N(n).
 \label{eq:hybrid-parity}
\end{equation}
Here the last equation follows from two successor computations and
the preceding equation.  Literal-equation inductions give
\begin{equation}
 N(2n)=0,\qquad N(2n+1)=1,\qquad N(3n)=N(n).
 \label{eq:hybrid-parity-scaling}
\end{equation}
For the first equation the successor step adds two to the argument of
$N$; the second follows by one computation.  In the third, increasing
$n$ adds three to the left argument and therefore applies $Z$ once to
both sides of the induction equation.

Finite computations now give
\[
 j(0)=1,\qquad j(1)=2,\qquad j(2)=0.
\]
The proof below needs just these three values and
$M(Sn)=j(M(n))$; it does not assume a division algorithm or a theorem
characterizing $M$ as remainder modulo three.

\subsection{A ramp with a global scaling identity}

Define
\[
 G(n)=I(n,0;u.u+M(u)+1).
\]
The central invariant is
\begin{equation}
 M(G(n))=N(n).
 \label{eq:hybrid-ramp-invariant}
\end{equation}
Its zero case is immediate.  For its successor step, retain the
single induction equation and split the Boolean value $N(n)$.
If it is $0$, the next $G$ state is $G(n)+1$; its $M$ value is
$j(0)=1$, equal to $N(Sn)$.  If it is $1$, the next $G$ state
is $G(n)+2$; two computations of $M$ change its value from $1$
to $2$ and then to $0$, again equal to $N(Sn)$.  This proves
the invariant by literal-equation induction.

It follows that
\[
 G(Sn)=G(n)+N(n)+1,\qquad G(n+2)=G(n)+3.
\]
Equation induction therefore yields
\begin{equation}
 G(2n)=3n.
 \label{eq:hybrid-ramp-double}
\end{equation}
A separate literal-equation induction proves
\begin{equation}
 G(3n)=3G(n)+N(n).
 \label{eq:hybrid-ramp-triple}
\end{equation}
For completeness, its successor calculation is
\begin{align*}
 G(3Sn)
 &=G(3n)+N(3n)+4
   =3G(n)+2N(n)+4,\\
 3G(Sn)+N(Sn)
 &=3G(n)+3N(n)+3+Z(N(n)).
\end{align*}
The expressions agree in both Boolean cases for $N(n)$.
The base is $0=0$.  In particular,
\eqref{eq:hybrid-ramp-triple} was proved internally; it was not
inferred from an intended floor-function interpretation of $G$.

Set
\[
 H(u)=C(N(u),G(u),u).
\]
The crucial choice is to retain $u$ when its parity is $1$.
The preceding equations give
\begin{equation}
 H(2u)=3u,\qquad H(1)=1,\qquad H(3u)=3H(u).
 \label{eq:hybrid-factor-map}
\end{equation}
To check the last identity, use $N(3u)=N(u)$ and split the Boolean
$N(u)$.  At $0$, both sides are $3G(u)$ by
\eqref{eq:hybrid-ramp-triple}; at $1$, both are $3u$.
This argument does not require that an arbitrary number has already
been represented as $2v$ or $2v+1$.

\subsection{A comparator from scalar iteration}

For every fixed update $h$, equation induction on $n$, keeping $m,a$
and all parameters fixed, proves the semigroup identity
\[
 I(m+n,a;h)=I(n,I(m,a;h);h).
\]
The base is the zero computation and the step applies $h$ to the
induction equation. This identity will also be used in the later
pairing construction.


For a fixed update $h$, equation induction on $y$, with its initializer
$a$ and all parameters of $h$ fixed, proves the rotation identity
\begin{equation}
 I(Sy,a;h)=I(y,h(a);h).
 \label{eq:hybrid-rotation}
\end{equation}
Both sides have base $h(a)$, and their successor equations apply the
same $h$ to the previous equality.  Here and below $I(y,a;h)$ is
an abbreviation for one fixed finite scalar term, not quantification
over functions.

By \eqref{eq:hybrid-factor-map}, another equation induction gives
\begin{equation}
 I(y,3a;H)=3I(y,a;H).
 \label{eq:hybrid-iterator-scaling}
\end{equation}
Its base is $3a=3a$.  In the step the fixed-parameter induction
equation rewrites the new left side as
\[
 H(3I(y,a;H))=3H(I(y,a;H)).
\]
The initializer parameter stays $a$ throughout this induction.
Equation induction also gives $I(y,1;H)=1$, since $H(1)=1$.

Define
\[
 E(x)=I(x,1;u.2u),\qquad
 L(x,y)=N\bigl(I(y,E(x);H)\bigr).
\]
Then the following entire recursion specification is provable:
\begin{equation}
 L(0,y)=1,\qquad L(Sx,0)=0,\qquad L(Sx,Sy)=L(x,y).
 \label{eq:hybrid-comparator}
\end{equation}
The first equation uses $E(0)=1$ and the fixed-point identity for
$H$.  The second uses $E(Sx)=2E(x)$ and $N(2a)=0$.
For the third,
\begin{align*}
 L(Sx,Sy)
 &=N\bigl(I(Sy,2E(x);H)\bigr)\\
 &=N\bigl(I(y,H(2E(x));H)\bigr)\\
 &=N\bigl(I(y,3E(x);H)\bigr)\\
 &=N\bigl(3I(y,E(x);H)\bigr)\\
 &=N\bigl(I(y,E(x);H)\bigr)=L(x,y).
\end{align*}
The successive steps use rotation, $H(2a)=3a$, iterator scaling,
and $N(3a)=N(a)$.  Thus $L$ has the comparison equations at all
arguments internally.  Its standard-model meaning as the indicator
of $x\le y$ is motivation, not a premise of the following argument.

\subsection{Saturated iteration gives predecessor}

We now need only the three equations
\eqref{eq:hybrid-comparator}.  Define
\[
 B_m(u)=C(L(Su,m),u,Su),\qquad
 W(t,m)=I(t,0;u.B_m(u)).
\]
The subscript in $B_m$ records a fixed parameter; $B_m$ is an open
term abbreviation, not a function variable.

Numerical cases on $e$ prove the selector equation
\[
 C(e,Sa,Sb)=S C(e,a,b).
\]
Equivalently one may use equation induction on $e$, whose induction
hypothesis is unused.  Therefore
\begin{align}
 B_{Sm}(Su)
 &=C(L(SSu,Sm),Su,SSu)\notag\\
 &=C(L(Su,m),Su,SSu)=S B_m(u),
 \label{eq:hybrid-conjugacy}\\
 B_{Sm}(0)&=C(L(S0,Sm),0,S0)=C(L(0,m),0,S0)=S0.
 \label{eq:hybrid-conjugacy-base}
\end{align}

Literal-equation induction on $t$, with $m$ fixed, proves
\begin{equation}
 W(St,Sm)=S W(t,m).
 \label{eq:hybrid-saturated-shift}
\end{equation}
The base is \eqref{eq:hybrid-conjugacy-base}.  The successor step is
\begin{align*}
 W(SSt,Sm)
 &=B_{Sm}(W(St,Sm))\\
 &=B_{Sm}(S W(t,m))\\
 &=S B_m(W(t,m))=S W(St,m).
\end{align*}
Only the second line uses the current induction hypothesis, at the
same fixed $m$.

Two subsequent equation inductions give
\begin{equation}
 L(Sn,n)=0,\qquad W(n,n)=n.
 \label{eq:hybrid-diagonals}
\end{equation}
The first uses the comparator recurrence.  The second uses
\eqref{eq:hybrid-saturated-shift} in its successor step.  The
substitution $t=m=n$ is made into that already proved universally
parameterized equation; no parameter of an active induction hypothesis
is changed.  Consequently
\begin{equation}
 W(Sn,n)=B_n(n)=C(L(Sn,n),n,Sn)=n.
 \label{eq:hybrid-extra-step}
\end{equation}

Finally define the finite scalar term
\begin{equation}
 P(m)=I\bigl(m,0;u.C(L(SSu,m),u,Su)\bigr).
 \label{eq:hybrid-predecessor-term}
\end{equation}
Its zero equation is immediate.  To justify its successor equation
without assuming a separate extensionality rule for iterator bodies,
put
\[
 V(t,m)=I\bigl(t,0;u.C(L(SSu,Sm),u,Su)\bigr).
\]
The comparator recurrence identifies this update pointwise with
$B_m(u)$.  Equation induction on $t$, with $m$ fixed, gives
$V(t,m)=W(t,m)$: their base values are $0$, and the successor step
uses the pointwise update equation and then the induction equation.
Now \eqref{eq:hybrid-extra-step} yields
\[
 \boxed{\qquad P(0)=0,\qquad P(Sm)=V(Sm,m)=W(Sm,m)=m.\qquad}
\]

\subsection{Dependency boundary}

All terms in this section are built from $0,S,I$ by finite substitution.
All inductions have a single equation as their invariant.  Numerical
cases and equality transport are the already reconstructed logical
resources.  Successor injectivity may now be obtained by applying $P$
to an assumed equation $Sx=Sy$.  Equality tests, reflection of
quantifier-free formulas, pairing, and the recovery of general
primitive recursion can then be developed in the established order.
They were not used in order to manufacture predecessor.

This argument is a written internal derivation in the stated theory.
The finite implementation checks do not constitute a proof-assistant
formalization of this entire metatheorem.

% Standalone section.  Requires only amsmath and amssymb in the parent document.
% No theorem environments, macros, or bibliography keys are assumed.
% All arithmetic claims in this section are proved here.

\section{From literal-equation induction to full intuitionistic PRA}
\label{sec:eic-completion}\label{app:arithmetic}

This section supplies an unconditional completion theorem. It is adapted
with attribution from the current forward-left draft~\cite{forward}. In the
scalar application used here, every function occurring in the proof is
already a scalar term: predecessor is supplied by
Appendix~\ref{sec:hybrid-bootstrap}, the selector is
$C(n,a,b)=I(n,a;u.b)$, addition is $x+y=I(y,x;u.Su)$, and
truncated subtraction is $d(x,y)=I(y,x;u.Pu)$. All other functions below
are finite compositions of these. Thus this proof does not presuppose
that arbitrary primitive recursion has already been compiled. The QFI
conclusion first applies to the scalar language; the general PR compiler
comes afterward.  Literal-equation
induction, together with the primitive-recursive defining equations,
predecessor coverage, and separation of $0$ and $1$, proves induction for every
quantifier-free formula.  In particular, the proof does not assume a faithful
numerical equality test, stability or decidability of equality, successor
uniqueness, or any commutation law for an iterator.  We construct the equality
test and prove its reflection property for arbitrary elements of the
numerical domain.

\subsection{The exact hypotheses}

Work in ordinary intuitionistic first-order logic with equality and genuine
falsity $\bot$.  The language has one numerical sort and a symbol for each
fixed, finite primitive-recursive description.  Distinct descriptions are
not identified merely because they compute the same standard function.
Let $\mathsf{PRE}$ contain the defining equations for zero, successor,
projections, composition, and primitive recursion.  In particular, choose
descriptions $P$ and $C$ with equations
\begin{align}
 P(0)&=0, & P(Sx)&=x, \label{eq:eic-predecessor}\\
 C(0,a,b)&=a, & C(Sx,a,b)&=b. \label{eq:eic-selector}
\end{align}
Write $1=S0$.  In addition to $\mathsf{PRE}$, assume
\begin{align}
 \mathsf{Cov}:&\qquad
       \forall x\,\bigl(x=0\lor x=S(Px)\bigr),
       \label{eq:eic-coverage}\\
 \mathsf{Sep}:&\qquad 0=1\longrightarrow\bot,
       \label{eq:eic-separation}
\end{align}
and induction for each \emph{literal equation} $E(n,\vec p)$:
\begin{equation}
 \forall\vec p\,
 \left[
   \left(E(0,\vec p)\land
       \forall n\bigl(E(n,\vec p)\longrightarrow E(Sn,\vec p)\bigr)
   \right)
   \longrightarrow\forall n\,E(n,\vec p)
 \right].
 \label{eq:eic-induction}
\end{equation}
Here the schematic expression $E$ is a single equality of numerical
terms; the displayed universal quantifiers range only over numbers.
Every description and every instance of the schema is finite external
syntax.

The axiom formulation \eqref{eq:eic-induction} and the usual contextual
induction rule are interchangeable.  For example, if
\[
 \Gamma\vdash E(0,\vec p),
 \qquad
 \Gamma,E(z,\vec p)\vdash E(Sz,\vec p),
\]
where $z$ is absent from $\Gamma$ and the passive parameters, discharge the
second premise's equation and universally introduce $z$.  The axiom then
gives $\Gamma\vdash\forall n E(n,\vec p)$, which can be instantiated at any
target term.  Conversely, assume the antecedent of the axiom and apply the
contextual rule with fresh variables.  This uses the ordinary intuitionistic
logical rules and permits arbitrary, including quantified, formulas in
$\Gamma$.

\subsection{Constructor separation and two elementary tests}

Define the sign description and the countdown description by
\begin{equation}
 \operatorname{sg}(u)=C(u,0,1),
 \qquad
 K(u)=C(u,1,Pu).
 \label{eq:eic-sign-countdown}
\end{equation}
Their defining equations imply
\begin{align}
 \operatorname{sg}(0)&=0,
 &\operatorname{sg}(Sx)&=1,\label{eq:eic-sign-computation}\\
 K(0)&=1,
 &K(Sx)&=x.\label{eq:eic-countdown-computation}
\end{align}

\paragraph{Constructor laws.}
Successor is injective: applying $P$ to $Sx=Sy$ gives $x=y$ by
\eqref{eq:eic-predecessor}.  Successor never equals zero.  Indeed,
$Sx=0$ gives $1=0$ by applying $\operatorname{sg}$ and using
\eqref{eq:eic-sign-computation}; symmetry and
\eqref{eq:eic-separation} give $\bot$.  Consequently $1\ne0$, where
$A\ne B$ abbreviates $A=B\to\bot$.  None of these arguments uses
induction.

\paragraph{The sign test.}
For every $u$,
\begin{align}
 &\operatorname{sg}(u)=0\lor\operatorname{sg}(u)=1,
       \label{eq:eic-sign-range}\\
 &\operatorname{sg}(u)=0\quad\longleftrightarrow\quad u=0,
       \label{eq:eic-sign-zero}\\
 &\operatorname{sg}(u)=1\quad\longleftrightarrow\quad u\ne0.
       \label{eq:eic-sign-one}
\end{align}
To prove the range assertion, split \eqref{eq:eic-coverage}.  The branch
$u=0$ gives sign $0$, and the branch $u=S(Pu)$ gives sign $1$.
The implication $u=0\to\operatorname{sg}(u)=0$ is computation.
For its converse, the successor branch of coverage would make
$\operatorname{sg}(u)=0$ imply $1=0$, a contradiction; in that branch
explosion gives $u=0$ as well.  Similarly, if
$\operatorname{sg}(u)=1$ and $u=0$, computation gives $0=1$.
Conversely, $u\ne0$ excludes the zero branch of coverage, and the
successor branch computes the sign as $1$.

Coverage and \eqref{eq:eic-selector} also give, for all $u,v$,
\begin{equation}
 C(u,v,v)=v.
 \label{eq:eic-selector-idempotence}
\end{equation}
Both coverage branches compute to $v$, so this identity again needs no
induction.

\paragraph{The countdown test and the notation $W$.}
For every $u$,
\begin{equation}
 K(u)=0\quad\longleftrightarrow\quad u=1.
 \label{eq:eic-countdown-zero}
\end{equation}
If $u=1=S0$, \eqref{eq:eic-countdown-computation} gives $K(u)=0$.
Conversely, assume $K(u)=0$ and split coverage.  If $u=0$, then
$K(u)=1$, contradicting separation, and explosion gives $u=1$.
If $u=S(Pu)$, computation gives $K(u)=Pu$; hence $Pu=0$ and
$u=S0=1$.

If a presentation uses a description named $W$ with defining equations
\[
 W(0)=1,\qquad W(Sx)=x,
\]
then it proves $K(u)=W(u)$ for every $u$ already from coverage.
In the zero branch both sides are $1$; in the branch $u=S(Pu)$ both
sides are $Pu$.  Thus $K$ and $W$ are provably equal on the whole domain,
even if they are distinct descriptions.  This equality does not appeal to
standard-model agreement or to uniqueness of an arbitrary recurrence.

\subsection{Four elementary uses of equation induction}

Choose descriptions for addition and truncated subtraction, with recursion
on the second argument:
\begin{align}
 x+0&=x, &x+Sy&=S(x+y),\label{eq:eic-addition}\\
 d(x,0)&=x, &d(x,Sy)&=P(d(x,y)).\label{eq:eic-subtraction}
\end{align}
Using these equations and \eqref{eq:eic-induction}, we now derive
\begin{align}
 0+y&=y,\label{eq:eic-left-zero}\\
 (Sx)+y&=S(x+y),\label{eq:eic-left-successor}\\
 d(Sx,Sy)&=d(x,y),\label{eq:eic-cancel-successors}\\
 d(Sa,a)&=1.\label{eq:eic-diagonal}
\end{align}
Each induction below has a single equation as its induction formula.

For \eqref{eq:eic-left-zero}, induct on $y$.  The base is $0+0=0$.
The step is $0+Sy=S(0+y)=Sy$.

For \eqref{eq:eic-left-successor}, hold $x$ fixed and induct on $y$.
At zero both sides are $Sx$.  At a successor the induction hypothesis
gives
\[
 (Sx)+Sy=S\bigl((Sx)+y\bigr)
       =S\bigl(S(x+y)\bigr)=S(x+Sy).
\]
Combining \eqref{eq:eic-left-successor} with the defining equation for
addition gives the useful rearrangement
\begin{equation}
 (Sy)+r=y+Sr.
 \label{eq:eic-shift-successor}
\end{equation}
Indeed, both sides are $S(y+r)$.

For \eqref{eq:eic-cancel-successors}, hold $x$ fixed and induct on $y$.
The base computes as
\[
 d(Sx,S0)=P(d(Sx,0))=P(Sx)=x=d(x,0).
\]
Applying $P$ to the induction equation gives the step:
\[
 d(Sx,S(Sy))=P(d(Sx,Sy))=P(d(x,y))=d(x,Sy).
\]

Finally, induct on $a$ for \eqref{eq:eic-diagonal}.  The base is
$d(S0,0)=S0=1$.  At a successor,
\[
 d(S(Sa),Sa)=d(Sa,a)=1,
\]
using \eqref{eq:eic-cancel-successors} and the induction hypothesis.

\subsection{An equation invariant that reflects actual equality}

The remaining arithmetic ingredient is the following equation, valid for
all $x,y$:
\begin{equation}
 C\bigl(\operatorname{sg}(d(x,y)),\ x,\ y+d(x,y)\bigr)=x.
 \label{eq:eic-residual-invariant}
\end{equation}
Its meaning is that a nonzero subtraction residual reconstructs the
original argument by addition.  The selector expresses this assertion as
one literal equation, so equation induction applies directly.

Hold $x$ fixed and induct on $y$.  At $y=0$, subtraction gives residual
$x$, and \eqref{eq:eic-left-zero} gives $0+x=x$.  The desired equation
is consequently $C(\operatorname{sg}(x),x,x)=x$, which is
\eqref{eq:eic-selector-idempotence}.

For the step, assume \eqref{eq:eic-residual-invariant} at $y$ and
abbreviate
\[
 r=d(x,y),\qquad r'=Pr=d(x,Sy).
\]
These are local term abbreviations, not additional function symbols.
Split coverage for $r'$.

If $r'=0$, then $\operatorname{sg}(r')=0$, and the successor instance
of \eqref{eq:eic-residual-invariant} reduces immediately to $x=x$.

In the other branch, $r'=S(Pr')$, so constructor separation gives
$r'\ne0$.  Split coverage for $r$.  Its zero branch would imply
$r'=P0=0$, a contradiction, from which the desired step equation follows.
In its successor branch,
\[
 r=S(Pr)=Sr',\qquad \operatorname{sg}(r)=1.
\]
The induction equation therefore computes to $y+r=x$.  Meanwhile,
\eqref{eq:eic-shift-successor} gives
\[
 (Sy)+r'=y+Sr'=y+r=x.
\]
Since $r'$ is a successor, $\operatorname{sg}(r')=1$, so this is
exactly the branch selected in the successor instance of
\eqref{eq:eic-residual-invariant}.  This completes the equation induction.
All case distinctions came from coverage; no induction on an implication
or a disjunction was used.

\paragraph{Reflection at residual $1$.}
We have, for arbitrary $a,b$,
\begin{equation}
 d(Sa,b)=1\quad\longleftrightarrow\quad a=b.
 \label{eq:eic-residual-reflection}
\end{equation}
For the forward direction, substitute $x=Sa$, $y=b$ into
\eqref{eq:eic-residual-invariant}.  If $d(Sa,b)=1$, the sign is $1$
and the selector returns its second branch, giving
\[
 b+1=Sa.
\]
The defining equation for addition gives $b+1=S(b+0)=Sb$.  Thus
$Sb=Sa$, and applying $P$ gives $b=a$, hence $a=b$.
Conversely, $a=b$ and \eqref{eq:eic-diagonal} give
$d(Sa,b)=d(Sa,a)=1$.

Define the equality characteristic by the fixed description
\begin{equation}
 e(a,b)=\operatorname{sg}\bigl(K(d(Sa,b))\bigr).
 \label{eq:eic-equality-characteristic}
\end{equation}
The sign, countdown, and residual reflection lemmas give
\begin{align}
 e(a,b)=0&\quad\longleftrightarrow\quad a=b,
       \label{eq:eic-equality-reflection}\\
 e(a,b)=0&\ \lor\ e(a,b)=1.
       \label{eq:eic-equality-range}
\end{align}
In detail, $e(a,b)=0$ is equivalent to $K(d(Sa,b))=0$ by
\eqref{eq:eic-sign-zero}, which is equivalent to $d(Sa,b)=1$ by
\eqref{eq:eic-countdown-zero}, and then to $a=b$ by
\eqref{eq:eic-residual-reflection}.  The range assertion is
\eqref{eq:eic-sign-range} applied to $K(d(Sa,b))$.

The reflection property concerns the logic's actual equality throughout
the domain.  It was not assumed as part of a numerical comparison symbol,
and its proof did not restrict $a$ or $b$ to standard numerals.
It also yields decidable equality.  In the $e(a,b)=0$ branch,
\eqref{eq:eic-equality-reflection} gives $a=b$.  In the $e(a,b)=1$
branch, an assumption $a=b$ would give $e(a,b)=0$ and therefore $0=1$,
contradicting separation.  Hence
\begin{equation}
 \forall a\,\forall b\,\bigl(a=b\lor(a=b\to\bot)\bigr)
 \label{eq:eic-equality-decidable}
\end{equation}
is a theorem, not an extra hypothesis.

\subsection{Compilation of all quantifier-free formulas}

For each fixed quantifier-free formula $A(\vec p)$, construct a term
$b_A(\vec p)$ by external recursion on its finite syntax.  Take equality,
$\bot$, $\land$, $\lor$, and $\to$ as the primitive formula constructors;
negation is implication into $\bot$.  Set
\begin{align}
 b_{s=t}&=e(s,t), & b_{\bot}&=1,\label{eq:eic-boolean-atoms}\\
 b_{A\land D}&=C(b_A,b_D,1),\label{eq:eic-boolean-and}\\
 b_{A\lor D}&=C(b_A,0,b_D),\label{eq:eic-boolean-or}\\
 b_{A\to D}&=C(b_A,b_D,0).\label{eq:eic-boolean-implies}
\end{align}
Every term here is a finite composition of fixed primitive-recursive
descriptions; if desired, projections and composition give it a single
description symbol.  No internal quantification over formulas or programs
is involved.

An external induction on the syntax of $A$ proves, within the arithmetic
theory,
\begin{equation}
 \bigl(b_A=0\lor b_A=1\bigr)
 \quad\text{and}\quad
 \bigl(b_A=0\longleftrightarrow A\bigr).
 \label{eq:eic-boolean-correctness}
\end{equation}
Here are the full connective checks.  In every case $b_A=1$ implies
$\neg A$: if $A$ held, the already established implication
$A\to b_A=0$ would contradict separation.

\paragraph{Atomic formulas and falsity.}
For an equation, both claims are
\eqref{eq:eic-equality-reflection}--\eqref{eq:eic-equality-range}.
For falsity the range assertion is $1=0\lor1=1$, with its right disjunct
proved by reflexivity.  Also $1=0\to\bot$ follows from separation and
symmetry, while $\bot\to1=0$ follows by explosion.  Thus $1=0$ is
equivalent to $\bot$, as required.

\paragraph{Conjunction.}
Split $b_A=0\lor b_A=1$.  In the zero branch, $A$ holds and
$b_{A\land D}=b_D$.  Its range follows from the range of $b_D$,
and its being zero is equivalent to $D$, hence to $A\land D$ under
the known $A$.  In the one branch, $\neg A$ holds and
$b_{A\land D}=1$.  Its range is immediate.  If it were zero,
separation would give $\bot$, hence $A\land D$; conversely an
assumption $A\land D$ contradicts $\neg A$, and explosion gives
$b_{A\land D}=0$.

\paragraph{Disjunction.}
In the branch $b_A=0$, the formula $A$ holds and
$b_{A\lor D}=0$.  The range is immediate, and both $b_{A\lor D}=0$
and $A\lor D$ hold, proving their equivalence.  In the branch
$b_A=1$, we have $\neg A$ and $b_{A\lor D}=b_D$.  Its range follows
from the range of $b_D$.  If its value is zero, then $D$ and hence
$A\lor D$.  Conversely, from $A\lor D$, the $A$ alternative
contradicts $\neg A$ and yields $b_D=0$ by explosion, while the
$D$ alternative yields $b_D=0$ directly.  This proves the required
equivalence in both alternatives.

\paragraph{Implication.}
In the branch $b_A=0$, the formula $A$ holds and
$b_{A\to D}=b_D$.  Its range follows from the range of $b_D$.
If its value is zero, then $D$, which gives $A\to D$.  Conversely,
$A\to D$ and the known $A$ give $D$, hence $b_D=0$.
In the branch $b_A=1$, we have $\neg A$ and
$b_{A\to D}=0$.  Its range is immediate.  To prove $A\to D$,
assume $A$; this contradicts $\neg A$, and explosion gives $D$.
Thus both $b_{A\to D}=0$ and $A\to D$ hold in this branch, and
their equivalence follows.  This argument uses the proved two-value
range of the finite characteristic.  It does not use excluded middle
for arbitrary formulas or change the surrounding intuitionistic logic.

This completes the external syntactic induction proving
\eqref{eq:eic-boolean-correctness} for every quantifier-free formula.

\subsection{Recovery of induction and the full theory}

Let $A(n,\vec p)$ be any fixed quantifier-free formula.  Suppose an
arbitrary retained context $\Gamma$ proves its base and its induction
step:
\[
 \Gamma\vdash A(0,\vec p),
 \qquad
 \Gamma,A(z,\vec p)\vdash A(Sz,\vec p),
\]
with the usual freshness condition on $z$.  By
\eqref{eq:eic-boolean-correctness}, the base gives
$\Gamma\vdash b_A(0,\vec p)=0$.  Under the additional assumption
$b_A(z,\vec p)=0$, the same equivalence gives $A(z,\vec p)$,
the given step gives $A(Sz,\vec p)$, and the equivalence gives
$b_A(Sz,\vec p)=0$.  Apply literal-equation induction to
$b_A(n,\vec p)=0$ and then apply its reflection implication.
The result is $\Gamma\vdash A(t,\vec p)$ for every target term $t$.
Thus the theory proves the full contextual quantifier-free induction
rule, including in arbitrary quantified contexts.

\paragraph{Completion theorem.}
In the common language of fixed primitive-recursive descriptions, the
theory
\[
 \mathsf{PRE}+\mathsf{Cov}+\mathsf{Sep}+\mathsf{EI}
\]
has exactly the same sequents as ordinary first-order intuitionistic
primitive recursive arithmetic with all quantifier-free induction
instances.  The forward simulation follows from the recovered
quantifier-free induction, successor injectivity, and successor separation;
all primitive-recursive defining equations and all logical rules are
shared.  For the reverse simulation, literal-equation induction is an
instance of quantifier-free induction, and successor separation gives
$\mathsf{Sep}$.  Quantifier-free induction proves $\mathsf{Cov}$:
at zero choose its left disjunct, and at a successor use
$Sx=S(P(Sx))$, which follows from $P(Sx)=x$, for its right disjunct.

Both interpretations therefore use the full domain, defined by $x=x$,
and preserve every function symbol, equality, logical connective, and
quantifier.  Both composites have the identity isomorphism, with graph
$x=y$.  The theories are consequently identity-domain bi-interpretable
and equiconsistent.  This completion theorem is unconditional once the
literal-equation induction rule has been obtained; it has no dependency
on a classical successor-uniqueness reduction.

\section{Pairing, primitive recursion, and internal numerical composites}
\label{sec:iterator-coding}
This appendix is adapted with attribution from the progress-closure
arithmetic development~\cite{progress}. Its prerequisites have now been
proved: scalar predecessor, arithmetic reflection, and full quantifier-free
induction. Pairing is a derived proof tool; it is not part of any checker.

\subsection{Pairing from a bounded square-root search}

Once QFI is available, a familiar square-based injection gives a
particularly direct compiler. Write $a\le b$ for $d(a,b)=0$.
The usual linear-order and semiring identities follow by fixed
quantifier-free inductions from the successor, addition, and monus
equations. Multiplication itself is the scalar term
\[
                        a\cdot b=I(b,0;u.u+a).
\]
Only primitive-recursive arithmetic is involved in the following
bounded root search; there is no minimization over an unbounded domain.

For a fixed parameter $n$, put
\[
 U_n(v)=C(d((Sv)^2,n),Sv,v),\qquad
 R_n(i)=I(i,0;v.U_n(v)),\qquad q(n)=R_n(n).
\]
The update increments $v$ exactly when $(v+1)^2\le n$, and
otherwise retains it. QFI in $i$, with $n$ fixed, proves
\[
 R_n(i)\le i,\qquad R_n(i)^2\le n,\qquad
 R_n(i)=i\ \vee\ n<(S(R_n(i)))^2.              \tag{Root invariant}
\]
The base has $R_n(0)=0$. If the update increments, its test
contradicts the strict-upper-bound alternative in the third invariant.
Thus the old state equals $i$, and the new state equals $Si$;
its square bound is the increment test. If the update does not increment, the
strict upper square bound holds and the state remains fixed.
An already established strict upper square bound forces every
subsequent update to remain fixed. These are exactly the cases
needed for the three displayed invariants.

At $i=n$, either the strict upper square bound already holds or
$R_n(n)=n$, in which case $n<(n+1)^2$ supplies it. Thus
\[
                        q(n)^2\le n<(q(n)+1)^2.           \tag{Root}
\]
Monotonicity of squares shows uniqueness of a number satisfying
Root. All formulas just used are quantifier free in the scalar
language, with fixed $n$ during the induction.

Now define
\[
 \pi(a,b)=(a+b)^2+a,\qquad
 L(n)=d(n,q(n)^2),\qquad H(n)=d(q(n),L(n)).
\]
For $k=a+b$, one has $k^2\le k^2+a<(k+1)^2$ since $a\le k$.
Root uniqueness gives $q(\pi(a,b))=a+b$, and therefore
\[
                    L(\pi(a,b))=a,\qquad H(\pi(a,b))=b.
\]
Surjectivity of $\pi$ is unnecessary. The original numerical domain
is never replaced by the range of this coding; pairing is used only
to store the finite state of a fixed primitive recursion. This
standard square pairing is a convenient reconstruction here, not
a claimed new arithmetic basis.

\subsection{Compiling each primitive recursion}

Suppose the fixed target description is
\[
 f(0,\vec p)=g(\vec p),\qquad
 f(Sn,\vec p)=h(n,f(n,\vec p),\vec p).
\]
Assume by external structural induction that the earlier $g,h$
already have scalar translations $g^*,h^*$. Define
\[
 W(u,\vec p)=\pi(SL(u),h^*(L(u),H(u),\vec p)),
\]
\[
 V(n,\vec p)=I(n,\pi(0,g^*(\vec p));u.W(u,\vec p)),
 \qquad f^*(n,\vec p)=H(V(n,\vec p)).
\]
The update has only its state variable bound, with every $p_i$
fixed. Equation induction proves $L(V(n,\vec p))=n$.
The two projection equations and scalar computation then prove
both defining equations for $f^*$ internally. Initial functions
and composition translate directly. This completes the translation
of every fixed finite PR description without adding a universal
evaluator.

\subsection{Why both numerical round trips are internal}

Two fixed scalar iterators with the same seed and provably equal
updates are internally equal at all counters: use literal-equation
induction comparing their values. Their base agrees, and their
successor equations plus congruence prove the step. The same
uniqueness argument works for ordinary primitive recursion.

For an ordinary PR description, induct externally on its finite
construction. The double translation satisfies its original
defining equations, so internal equation induction identifies it
with the original operation. For a native scalar description,
earlier translated updates agree by the external induction
hypothesis; scalar uniqueness identifies the double-translated
iterator. These proofs hold at arbitrary internal counters.

Thus the domain relation $x=y$ is provably a total bijection
preserving every fixed numerical operation in both composites.
The formula constructions and round trips in
Section~\ref{sec:choice-introduction} give the corresponding formula
composites. Both interpretations preserve falsity, taking native
$0=1$ to an ordinary contradiction and ordinary $\bot$ to native
$0=1$. Finite proof translations therefore establish equiconsistency.

% Standalone mathematical module; requires amsmath, amssymb, amsthm, url.
% The surrounding document should supply theorem, lemma, and corollary environments.
% Written derivation, not a proof-assistant verification.
\section{Removing literal zero from the source of transport}
\label{zs:section}

The following is a restriction of a mathematical inference rule, rather
than a change in the spelling of zero. Throughout this section,
``nonzero source'' means \emph{not syntactically the single term $0$}.
Every variable and every term headed by $S$ or $I$ is a permitted source,
including terms whose numerical value is zero.

Let $\mathcal C$ be the native scalar calculus with terms
\[
 t::=x\mid0\mid St\mid I(n,a;u.h),
\]
the two computation schemas
\[
 I(0,a;u.h)=a,\qquad
 I(Sn,a;u.h)=h[I(n,a;u.h)/u],
\]
formulas generated by equality and the $Q,D$ binders, and the Hyp, QE,
DE, case-DI and equation-progress-DI rules stated in the main inventory.
The same argument applies to the original presentation with $D$ object
formulas and $Q$ resources in contexts. Let $\mathcal C^{-}$ replace
unrestricted transport by
\[
 \frac{\Gamma\vdash s=t\qquad\Gamma\vdash F(s)}
      {\Gamma\vdash F(t)}
 \quad\mathsf{Tr}^{-},\qquad s\not\equiv0.
 \tag{Tr$^-$}\label{zs:restricted-rule}
\]
Here $F$ is an arbitrary native numerical template, and $s\not\equiv0$
is a condition on finite term syntax. All other rules are unchanged.

\begin{theorem}[Literal-zero source elimination]
\label{zs:main}
For every native context $\Gamma$ and formula $A$,
\[
 \Gamma\vdash_{\mathcal C^{-}}A
 \quad\Longleftrightarrow\quad
 \Gamma\vdash_{\mathcal C}A.
\]
In particular, the ordinary whole-domain realization of
$\mathcal C^{-}$ is bi-interpretable and equiconsistent with full
first-order iPRA through the same interpretations as the native scalar
core. No change of numerical domain or interpreted equality is needed.
\end{theorem}

We prove the theorem without assuming arbitrary numerical substitution
admissibility for $\mathcal C^{-}$: such substitution can turn a variable
source into the forbidden literal zero. Every occurrence of a helper
equation at an arbitrary term is either obtained directly from its
displayed derived rule at that term, or from a derivation whose transport
sources retain an $S$ or $I$ head under the indicated substitution.

\subsection{Structural operations, cases, and equation induction}

Weakening, ordinary-assumption cut, and capture-free renaming are
admissible by induction on derivations. These operations do not turn a
nonzero source into literal zero. As usual, fresh parameters are renamed
before a cut. We also use the following restricted substitution fact:
a capture-free substitution is admissible for a particular finite
derivation if it leaves every primitive transport source in that
derivation syntactically different from zero. This fact follows by the
same induction; it is not unrestricted substitution admissibility.

\begin{lemma}[Cases and equation induction need no transport]
\label{zs:cases-induction}
In $\mathcal C^{-}$ the following rules are derivable:
\[
 \frac{\Gamma\vdash A(0)\qquad\Gamma\vdash A(Sz)}
      {\Gamma\vdash A(t)}\quad\mathsf{NC},
 \qquad
 \frac{\Gamma\vdash E(0)\qquad\Gamma,E(z)\vdash E(Sz)}
      {\Gamma\vdash E(t)}\quad\mathsf{EI}.
\]
In NC, $A$ is any native formula; in EI, $E$ is one literal equation.
The parameter $z$ is fresh for $\Gamma$, $t$, and the fixed parameters
of the matrix. Other parameters stay fixed throughout EI.
\end{lemma}
\begin{proof}
Choose an unused witness coordinate $x$. In NC, weaken the successor
proof by $A(0)$ and apply case-DI, with witness zero, to obtain
$D_{x,y}A(y)$. DE opens $Q_yA(y)$, because $x$ is unused. QE with the
given base proof and target $t$ yields $A(t)$, and DE closes.
For EI use progress-DI instead, with the given equation step, and use
the same DE/QE construction. No transport rule occurs.
\end{proof}

These rules will be applied to unconditional helper equations before
they are weakened into the context containing $0=t$. We never use NC
with a case parameter occurring in its retained context.

\subsection{The equality operations available immediately}

\begin{lemma}[Reflexivity and safe equality operations]
\label{zs:safe-equality}
Reflexivity $r=r$ is derivable for every term $r$. From $p=q$ with
$p\not\equiv0$, one can derive $q=p$ and $H(p)=H(q)$ for any numerical
template $H$. Every direct computation can be used in either direction
as an equation; using the reversed equation as a transport major is
permitted only when its new source is syntactically nonzero.
\end{lemma}
\begin{proof}
Put $b=I(0,r;u.0)$. The computation $b=r$, used as both major and
minor of transport with template $v=r$, gives $r=r$. Its source is
the nonempty iterator term $b$.
For symmetry of $p=q$, transport the reflexivity $p=p$ through the
template $v=p$. For numerical congruence, transport the reflexivity
$H(p)=H(p)$ through $H(p)=H(v)$. Both use the permitted source $p$.
Every computation equation has an $I$-headed left side, so its
symmetry follows by the same construction.
\end{proof}

An occurrence-selective template may keep some occurrences of a term
fixed while abstracting other occurrences. This is ordinary
capture-free template substitution, not an extra inference rule.

\begin{lemma}[Joining two nonzero terms known to be zero]
\label{zs:zero-join}
Suppose $p,q\not\equiv0$ and $\Gamma$ proves $p=0$ and $q=0$.
Then $\Gamma\vdash p=q$ using only $\mathsf{Tr}^{-}$ and computation.
\end{lemma}
\begin{proof}
Let
\[
 z_0=I(0,0;u.0),\qquad b_p=I(0,p;u.0).
\]
Computation gives $b_p=p$. Start with $b_p=b_p$ and use the major
$p=0$ in the template $I(0,v;u.0)=b_p$. This gives $z_0=b_p$.
Symmetry at the nonzero source $z_0$ gives $b_p=z_0$.
Now transport through $v=z_0$ using the computation $b_p=p$, whose
source is $b_p$, to obtain $p=z_0$.
The identical construction gives $q=z_0$. Symmetry at source $q$
gives $z_0=q$, and transport at source $z_0$ in $p=z_0$ gives $p=q$.
The sources used are $p,q,z_0,b_p,b_q$; all are syntactically nonzero.
No reverse transport from zero has been used.
\end{proof}

The zero-join lemma is deliberately placed before the next construction.
The positive case of its singleton test needs precisely this lemma;
ordinary transitivity through zero is not yet available.

\subsection{A singleton-one test without predecessor}

Define the finite scalar terms
\[
\begin{aligned}
 C(v,a,b)&:=I(v,a;u.b),& Z(v)&:=C(v,1,0),\\
 K(v)&:=C(v,0,1),& H(v)&:=C(v,1,2),\\
 N(n)&:=I(n,0;u.Z(u)),& R(n)&:=I(n,0;u.H(u)),\\
 L(n)&:=N(R(n)).
\end{aligned}
\tag{Tests}\label{zs:tests}
\]
The update variable in $C$ is unused. Here $1=S0$ and $2=S1$.
The standard-model description of $R$ is $0,1,2,2,\ldots$; the proof
below establishes the needed equations internally and does not assume
this description.

\begin{lemma}[Uniform test equations]
\label{zs:singleton}
The following equations have derivations in $\mathcal C^{-}$,
uniformly in the displayed argument:
\[
 H(v)=S K(v),\qquad R(1)=1,\qquad R(SSv)=2,
\]
\[
 N(1)=1,\quad N(2)=0,\quad
 L(1)=1,\quad L(SSv)=0,\quad
 \boxed{L(Sv)=Z(v)}.
 \tag{Singleton}\label{zs:singleton-equations}
\]
All transport sources in these helper derivations have an $I$ or $S$
head, except for the explicit nonzero endpoints in the zero-join
lemma; those endpoints also have an $I$ head in this application.
\end{lemma}
\begin{proof}
For $H(v)=S K(v)$ use NC. At zero both sides reduce to $1$; at a
successor both reduce to $2$. To join two equations $p=j$, $q=j$
with $j=1$ or $j=2$, reverse $q=j$ at its nonzero source $q$, then
transport at the nonzero numeral $j$. Each reduction of these
particular terms is a forward replacement of an iterator computation.
This proves the helper at an arbitrary target $v$ by NC, without
substituting zero into a variable-source transport.

The equations $N(1)=1$ and $N(2)=0$ are finite computations:
$N(0)=0$, $N(1)=Z(N(0))=1$, and $N(2)=Z(N(1))=0$.
Every replacement source in these calculations is an iterator term.
Likewise $R(1)=H(R(0))=H(0)=1$.

Computation and the first helper give, for any $v$,
\[
 R(Sv)=H(R(v))=S K(R(v)).
\]
Use this equation at the nonzero source $R(Sv)$ in
$R(SSv)=H(R(Sv))$. The resulting right side is
$H(S K(R(v)))=2$ by the successor computation of $C$.
Thus $R(SSv)=2$. These substitutions are all at $I$-headed sources.
Substitute $R(1)=1$ and $R(SSv)=2$ inside $N$ and use the two
computed values of $N$ to obtain $L(1)=1$ and $L(SSv)=0$.

Finally apply NC to the matrix $L(Sv)=Z(v)$.
At zero, $L(1)=1$ and $Z(0)=1$ are joined through the nonzero numeral
$1$. At a successor, $L(SSv)=0$ and $Z(Sv)=0$ are joined by
Lemma~\ref{zs:zero-join}; both left sides are nonempty iterator terms.
This proves the boxed equation at every target. Its use is therefore
an instance of a displayed derived rule, not an appeal to unrestricted
substitution admissibility.
\end{proof}

\begin{lemma}[An unconditional zero-filter identity]
\label{zs:filter}
For every term $t$,
\[
                 t=C(Z(t),t,0)
                 \tag{Filter}\label{zs:filter-equation}
\]
is derivable in the empty context of $\mathcal C^{-}$.
\end{lemma}
\begin{proof}
Apply NC to $v=C(Z(v),v,0)$, with the target $t$.
At zero, forward computations give
$C(Z(0),0,0)=C(1,0,0)=0$; reverse this equation at its nonzero
left side. At $Sv$, forward computations give
$C(Z(Sv),Sv,0)=C(0,Sv,0)=Sv$; again reverse at the nonzero
iterator left side. Thus the base and successor proofs have no
retained assumption depending on the case variable. All numerical
arguments of the final NC target are permitted. In particular, this
identity is proved before any assumption $0=t$ is used.
\end{proof}

\subsection{Turning a zero-source equality around}

\begin{lemma}[Zero-source symmetry]
\label{zs:zero-symmetry}
For every term $t$,
\[
                     0=t\ \vdash_{\mathcal C^{-}}\ t=0.
                     \tag{ZS}\label{zs:zero-symmetry-equation}
\]
\end{lemma}
\begin{proof}
If $t\equiv0$, use reflexivity. Otherwise $t$ is a permitted source.
Work in the fixed context $\Gamma=\{0=t\}$ and put
\[
                  A(n):=I(n,t;u.Su).
\]
We first prove $n=A(n)$ by EI, keeping $t$ fixed.

For the base, computation gives $A(0)=t$. Its symmetry gives
$t=A(0)$. Since $t\not\equiv0$, transport at source $t$ in the
retained assumption $0=t$ gives $0=A(0)$.

For the step, retain $z=A(z)$ with $z$ fresh. This source is a
variable, hence is permitted. Numerical successor congruence from
Lemma~\ref{zs:safe-equality} gives $Sz=S A(z)$.
Computation gives $A(Sz)=S A(z)$, and its symmetry gives
$S A(z)=A(Sz)$. Its source $S A(z)$ is also permitted, so transport
in $Sz=S A(z)$ gives $Sz=A(Sz)$. This is exactly the EI step.
The induction hypothesis is used at the same fixed parameter $t$.

At target $1$, EI gives $1=A(1)$. Forward computation gives
$A(1)=St$: the successor equation produces $S A(0)$, and the
nonzero-source computation $A(0)=t$ is used inside $S$.
Transport at source $A(1)$ consequently gives
\[
                              1=St.
\]
Its source $1$ is permitted. Numerical congruence for $L$ gives
$L(1)=L(St)$. Use $L(1)=1$ at source $L(1)$, and
$L(St)=Z(t)$ at source $L(St)$, obtaining $1=Z(t)$.
Symmetry at source $1$ gives $Z(t)=1$.

Weaken the unconditional Filter identity into $\Gamma$.
Transport at the iterator source $Z(t)$ changes its right side to
$C(1,t,0)$. The latter computes to zero, and its source is again
an iterator term. The resulting equation is $t=0$, as required.
\end{proof}

This proof uses cases only for unconditional helper equations. The
sole induction with the retained assumption $0=t$ has the literal
equation $n=A(n)$ as invariant. It does not perform a forbidden
case distinction while silently changing the retained context.

\subsection{Recovering equality algebra and formula transport}

\begin{lemma}[Full numerical equality algebra]
\label{zs:full-equality}
The calculus $\mathcal C^{-}$ admits unrestricted symmetry,
numerical congruence, and transitivity. In particular it admits
root replacement
\[
 \frac{\Gamma\vdash s=t\qquad\Gamma\vdash s=r}
      {\Gamma\vdash t=r}.
\]
\end{lemma}
\begin{proof}
Symmetry uses Lemma~\ref{zs:safe-equality} when the source is not
zero and Lemma~\ref{zs:zero-symmetry} otherwise.

For numerical congruence from $s=t$, a nonzero $s$ is handled by
Lemma~\ref{zs:safe-equality}. If $s\equiv t\equiv0$, use
reflexivity. If $s\equiv0$ and $t\not\equiv0$, first derive $t=0$,
apply permitted numerical congruence at source $t$ to obtain
$H(t)=H(0)$, and use the now established full symmetry.

For transitivity from $r=s$ and $s=t$, a nonzero middle term $s$
is handled by one permitted transport. If $s\equiv0$ and either
endpoint is literally zero, a given premise is already the
conclusion. Otherwise $r,t\not\equiv0$; reverse $0=t$ to get
$t=0$ and apply the zero-join lemma to $r=0$ and $t=0$.
Root replacement follows by symmetry and transitivity.
These constructions are acyclic: zero-source symmetry was proved
before any unrestricted numerical congruence or transitivity.
\end{proof}

\begin{lemma}[Full transport]
\label{zs:full-transport}
For every native numerical template $F(v)$, the rule
\[
 \frac{\Gamma\vdash s=t\qquad\Gamma\vdash F(s)}
      {\Gamma\vdash F(t)}
\]
is admissible in $\mathcal C^{-}$, without a source restriction.
\end{lemma}
\begin{proof}
Use external induction on the finite number of formula binders in
$F$, uniformly in retained contexts and fixed numerical parameters.
For an equation $F(v)\equiv A(v)=B(v)$, use unrestricted numerical
congruence to obtain $A(s)=A(t)$ and $B(s)=B(t)$. Join these with
$A(s)=B(s)$ using the established equality algebra.

For $F(v)=D_{x,y}B(v,x,y)$, eliminate $F(s)$ at a fresh $a$,
retaining $Q_yB(s,a,y)$. Introduce the desired $F(t)$ by case-DI
with witness $a$. Its local assumption is $B(t,a,0)$.
Full symmetry gives $t=s$, and the external induction hypothesis
for the smaller matrix transports the assumption to $B(s,a,0)$.
QE gives $B(s,a,Sz)$, and the induction hypothesis in the forward
direction gives $B(t,a,Sz)$. Close case-DI and DE. All new
parameters can be chosen fresh for the source equality and the
retained context.

If $Q$ is first-class formula syntax, first observe that its case
introduction is derived without transport: introduce
$D_{u,y}B(y)$ with unused $u$ by case-DI and eliminate it with the
one-node continuation $Q_yB(y)$. For $F(v)=Q_yB(v,y)$, use this
derived introduction for the target. Convert its base assumption
backwards by the induction hypothesis, apply the old Q resource,
and convert the successor conclusion forwards. This is the same
smaller-matrix argument and uses no new numerical induction.
\end{proof}

\begin{proof}[Proof of Theorem~\ref{zs:main}]
Every restricted derivation is an unrestricted one. Conversely,
replace each unrestricted transport node by the derivation supplied
by Lemma~\ref{zs:full-transport}, recursively translating its
premises first. Other inference nodes are unchanged. The result has
the same context and conclusion. Thus the full native consequence
relation, not just closed theorems, is preserved. The existing
native/ordinary and ordinary/iPRA interpretations and both identity
composites therefore transfer unchanged.
\end{proof}

\subsection{Compatibility with restricted computation certificates}
\label{zs:certificate-computation}

Some compact checkers require the initializer of a computation packet
to have nonempty syntax, and encode a successor counter by a nonempty
predecessor. These are certificate restrictions, not restrictions
on the two mathematical computation schemas. They do not introduce
a circular dependence on zero-source transport in the preceding proof.

Use the inherited certified-zero term
\[
                         c_0=I(2,1;u.0).
\]
Its equation $c_0=0$ is a directly admitted successor computation:
the predecessor and initializer are both $1$. For an arbitrary
computation instance with literal-zero initializer or literal-zero
successor predecessor, form its whole equation as a template $J(v)$,
replacing just these offending occurrences by $v$. The instance
$J(c_0)$ is a directly admitted computation. Transport through this
equation template using $c_0=0$ yields $J(0)$, the desired original
instance. Both computation premises are admitted directly, and the
only transport source is the nonempty term $c_0$.

This is exactly the inherited three-node normalization in
\path{Research/certified_zero/tests/producer.py}, function
\texttt{expand\_computation}: one transport and two computation
premises. Its source is $c_0$, even when both defects are corrected
simultaneously. Passive zero occurrences in an update are left alone.
Thus all computation instances used in the mathematical proof can
be expanded before compiling to the source-restricted checker.

\subsection{Scope of the result}

The new statement is the admissibility of all omitted literal-zero
source transports. The native term/formula language, arithmetic
schemas, and target theory remain the inherited ones. This theorem
justifies replacing a two-way equality/Q major-pattern test by a
source-discriminated test in a suitable certificate protocol. It
does not assert a minimum proof-checker size, a shortest proof
expansion, or a new characterization of the primitive recursive
functions. The displayed proof is a written universal derivation;
executable certificates, when supplied, validate particular
instances rather than substitute for this argument.

\section{A uniform arity encoding and its checking theorem}
\label{sec:arity-checker}

The arithmetic and object logic remain the scalar-iteration Eq/$\D$
calculus specified above.  The change is in its finite certificate
representation.  A paired formula is unary, an equation is binary, and
an internal closure is ternary:
\[
 \D_{x,y}B\longleftrightarrow[B],\qquad
 s=t\longleftrightarrow[s,t],\qquad
 \Q_yB\longleftrightarrow[B,B,B].
 \tag{Arity encoding}\label{eq:arity-encoding}
\]
The three copies in a closure are representation redundancy: only the
third has mathematical content.  Their benefit is that tripling a unary
major constructs exactly its closure.  Tripling any other nonempty
formula array produces an array of arity at least six, which will denote
truth.  This makes the other major cases neutral without a separate
test.

The previous generated-assumption and constructed-computation interfaces
are inherited from the corpus~\cite{generated,constructed}.  The new
encoding makes a single arithmetic expression, $5-\operatorname{arity}$,
perform the necessary depth adjustments.  Its soundness depends on an
explicit interpretation of surplus syntax.  It would not be sound to
drop arity checks while continuing to treat every surplus array as an
unconstrained new connective.

\subsection{The complete sorted interpretation of arrays}
\label{subsec:sorted-arrays}

Write $\top$ for the fixed equation $0=0$.  A protected numerical slot
has width two.  The numerical root has depth three: the coordinates
$0,1,2$ all denote zero, and under $p$ protected slots the depth is
$3+2p$.  Coordinates $3+2j$ and $4+2j$ name slot $j$, counted from
the innermost slot, for $0\leq j<p$.  Formula depth is $2p$ under the
same $p$ slots.  Negative coordinates admitted by the current global
floor name rigid numerical parameters.

The decoder is sorted by the position in which a tree is used.  Let
$\rho$ be its ordered environment of protected numerical slots.
The notation $(a,\rho)$ means prepend a slot of value $a$.
A superscript $\sharp$ denotes a child decoded in the indicated
environment.

\begin{center}
\begin{tabular}{cll}
\toprule
Numerical arity & Meaning & Relevant child environment\\
\midrule
$0$ & $0$ & No child\\
$1$: $[b]$ & $Sb^\sharp$ & $(0,0,\rho)$\\
$2$ & $0$ & Annotations ignored\\
$3$: $[n,a,h]$ & $I(n^\sharp,a^\sharp;v.h^\sharp)$ & $(0,\rho)$ for $n,a$; $(v,\rho)$ for $h$\\
$\geq4$ & $0$ & All children ignored\\
\bottomrule
\end{tabular}
\end{center}

\begin{center}
\begin{tabular}{cll}
\toprule
Formula arity & Meaning & Relevant child environment\\
\midrule
$0$ & $\top$ & No child\\
$1$: $[B]$ & $\D_{x,y}B^\sharp$ & $(y,x,\rho)$\\
$2$: $[s,t]$ & $s^\sharp=t^\sharp$ & $\rho$ numerically\\
$3$: $[P,U,B]$ & $\Q_yB^\sharp$ & $(y,\rho)$ for $B$\\
$\geq4$ & $\top$ & All children ignored\\
\bottomrule
\end{tabular}
\end{center}

The first two entries of a ternary formula are annotations.  A canonical
closure copies its body into both of them, as in
\eqref{eq:arity-encoding}.  Successfully admitted scalar formula padding
also denotes $\top$.  An internal closure has the ordinary meaning
$B(0)\to\forall yB(y)$; it is not an equation and is not eligible for
compound induction.

Canonical encoded object expressions use arrays and integer coordinates,
respect these slots, and contain no high-arity node.  The checker itself
still generates high-arity truth entries in administrative contexts,
even during a canonical proof.  In particular, constructing a
successor beneath positive bound coordinates requires passing through
two dummy zero slots.  Constructing an iterator requires one new slot:
zero in its counter and seed, and its genuine state variable in the
update.  Rigid negative parameter names are unaffected by these shifts.

The additional raw shapes are certificate aliases with the displayed
meanings.  They are not independent predicate or function symbols added
to the object theory.  Our exact completeness statement concerns the
canonical native language.  We do not require every possible raw spelling
of a semantic alias to be interchangeable under the same raw certificate.

\subsection{One depth calculation replaces the lookup table}

For every array the opener recursively adds $5-\operatorname{arity}$
to its depth.  The meaningful cases are
\[
\begin{array}{c|ccc}
 \operatorname{arity}&1&2&3\\\hline
 \text{depth increment}&4&3&2.
\end{array}
\]
Unary $\D$ adds two numerical binders; unary successor adds two dummy
zero slots.  Binary equality changes formula depth $2p$ to numerical
depth $3+2p$.  Ternary $\Q$ adds one numerical binder; ternary iteration
adds one state slot, or one dummy slot in its counter and initializer.
The same calculation therefore covers all relevant constructors.

Binary numerical annotations denote zero, so their parity-changing
traversal has no mathematical effect.  Empty arrays have no children.
At arity at least four, the entire node already denotes a constant.
Even a zero or negative depth increment below such a node cannot expose
one of its ignored children as a meaningful variable.  Traversal may
still reject an alias; it cannot change its stipulated interpretation.

This is a semantic assertion, not raw idempotence of the opener.
For example, at floor $-1$, opening a six-element array of zeros with
image $-1$ can yield six copies of $-1$, while a later no-image opening
of that result fails at a negative child depth.  Both formula arrays
denote $\top$.  The proof below only concerns successful calls at meaningful
sorted positions.  Canonical input object expressions contain no surplus
arrays; the administrative exports that generate them are justified
separately as truth entries.

\subsection{Opening and scope}
\label{subsec:opening-lemma}

Write $i(a,t,d)$ for the executable opener at floor $k$.  It preserves a
scalar when $k\leq a<d$.  At a numerical position, every other relevant
scalar is the distinguished substitution hole; the supplied image is
checked at numerical root depth three before insertion.  At a formula
position scalar padding is truth and carries no numerical hole.
Missing images at used numerical holes reject.

\begin{lemma}[Semantic opening]
\label{lem:semantic-opening}
For every successful call at a meaningful formula root of depth $2p$
or numerical root of depth $3+2p$, with $p\geq0$, decoding its result
is capture-free filling of the distinguished hole in the decoded input
template.  These depths protect $p$ numerical slots; admitted rigid
parameters are preserved.  If an image is semantically unused, zero
may be chosen as its mathematical representative.  Recursive calls
at negative depths inside ignored annotations need no independent
sorted interpretation.
\end{lemma}

\begin{proof}
Induct on the finite raw tree.  At a relevant scalar, the interval
comparison is exactly the protected-coordinate test.  At a meaningful
unary, binary, or ternary node, the preceding tables supply precisely
the required additional slots.  At numerical binary annotations,
high-arity nodes, and ignored closure fields, the relevant mathematical
meaning is constant or absent, so the conclusion is immediate
regardless of traversal inside the annotation.

An inserted image is checked at its own numerical root.  Its free
parameters are rigid negative names, and it has no unbound positive
coordinate there.  Insertion beneath protected slots therefore captures
none of its parameters; its internal binders retain their relative
indices.  This proves the inductive step for substitution as well as
for mere traversal.
\end{proof}

A successful no-image check consequently preserves the interpretation
of its checked closed expression.  Raw structural equality of checked
outputs implies equality of their decoded meanings.  The checker does
not infer equality merely from a semantic comparison of two different
programs or numerical terms.

The global floor decreases by one at each non-hypothesis continuation.
If its old value is $K$, the integer $K-1$ is the designated fresh
eigenparameter.  With integer coordinates it is the only new name.
If fractional coordinates are accepted, other names in $[K-1,K)$ may
also occur in the finite certificate.  These are additional fresh
parameters.  They must not be silently equated to the designated
integer eigenparameter.

Previously checked relevant goals and retained hypotheses mention none
of these new names: their earlier no-image checks would reject a
relevant numerical occurrence below the old floor.  Coordinates inside
ignored annotations do not contribute free variables to the decoded
formula.  This is the freshness invariant used in every rule below.

\subsection{A computation hypothesis only needs a zero instance}
\label{subsec:zero-supply}

The shared continuation can open a generated arithmetic equation at a
fresh parameter.  That equation need not be a theorem uniformly in the
new parameter.  The following admissible rule is sufficient.

\begin{lemma}[Supply with a known zero instance]
\label{lem:zero-supply}
Suppose $\vec a$ is fresh for $\Gamma,F$, the calculus proves
$H(\vec0)$, and $\Gamma,H(\vec a)\vdash F$.  Then
$\Gamma\vdash F$.
\end{lemma}

\begin{proof}
Substitute zero for every member of $\vec a$ in the second derivation.
Freshness leaves $\Gamma,F$ unchanged.  Cut in the derivation of
$H(\vec0)$.  Both substitution and cut are admissible in the native
calculus.  This is a finite proof transformation and does not use a
consistency assumption.
\end{proof}

All names newly admitted when the floor changes belong to $\vec a$;
for noninteger aliases this includes the finite collection in
$[K-1,K)$ that actually occurs.  On decoding, their simultaneous zero
substitution agrees with opening the generated expression at zero using
the old floor.  Previously admitted parameters are preserved.  New
rigid names and unprotected placeholders both become zero.  The opener
may keep different raw zero spellings, but their decoded values agree.

The constructed-computation packet supplies a triple $s=[n,b,r]$ and a
mode $u$.  In successor mode its generated equation has raw form
\[
 \bigl[\,[[n],b,r],\ i(r,s,3)\,\bigr],
 \tag{Step supply}\label{eq:step-supply}
\]
while in zero mode it has form
\[
                    \bigl[\,[u,b,r],\ b\,\bigr].
 \tag{Zero supply}\label{eq:base-supply}
\]
Here successor mode means that $u$ is truthy; the zero mode has falsy
$u$.  These expressions describe actual arrays, before their final
opening.  In \eqref{eq:step-supply} the counter is constructed as a
successor, rather than accepted on an unsupported assertion that it is
one.  This construction is inherited from the corresponding draft
interface~\cite{constructed}.

\begin{lemma}[Sound constructed computation]
\label{lem:constructed-supply}
Every successful arithmetic supply has a decoded zero instance
derivable from the two scalar-iteration computation schemas.
\end{lemma}

\begin{proof}
First apply the simultaneous zero substitution described above.  In
zero mode every successful falsy numerical counter denotes zero, so
\eqref{eq:base-supply} is an $I0$ instance.  In the exact-list profile
the relevant accepted aliases are an empty array, numeric zero, false,
and numerically equal zero values.  Null, strings, and dictionaries
cannot enter a meaningful numerical position through an unchecked
coercion.

In successor mode the left counter is an actual successor.  If the
update uses its state in a semantically relevant position, opening it
invokes the no-image check of the recursive source $s$ at numerical
root three.  This checks the entire source, including its update.
Only coordinates three and four are aliases for its single state slot.
A genuinely escaping coordinate in that update is revisited during
the recursive-image check and rejects.  Its counter and initializer
are likewise closed with their prescribed dummy zero slot.

Inserting the extra dummy slots of the displayed successor does not
alter such a checked predecessor: extra unbound slots have value zero,
and actual free parameters are rigid.  Lemma~\ref{lem:semantic-opening}
therefore identifies the right side of \eqref{eq:step-supply} with the
update applied to exactly the recursive predecessor value occurring
on its left.  This is the $IS$ schema.

If the update has no semantically relevant state occurrence, its
decoded value is constant in that variable.  Its zero-normalized value
is the same in the source update and on the right, independently of
the unused recursive image; this is again an $IS$ instance.  Traversal
inside an ignored annotation may still inspect the image and reject.
That removes a certificate and does not invalidate the argument for a
successful one.

Finally, successful three-field source extraction cannot hide another
JSON operator.  A scalar is not iterable.  A three-character string,
or a three-key JSON dictionary, necessarily puts a nonempty string in
a relevant source field, which the inspected opener rejects.  The
successful source is therefore the stated array triple.
\end{proof}

The recursive-image check in the state-dependent case is essential.
Merely opening a final displayed equation at zero would not establish
that its two state interpretations agree.  Combining
Lemma~\ref{lem:constructed-supply} with Lemma~\ref{lem:zero-supply}
justifies both actual image choices of the common continuation, including
the branch that retains a fresh-parameter hypothesis.

\subsection{Audit of the logical continuations}
\label{subsec:logical-audit}

\paragraph{Hypothesis closure.}
The initial opener check and membership comparison select the same
decoded context formula.  Additional aliases can make a spelling less
convenient, but cannot identify an equation with a false denotation
solely by semantic evaluation.

\paragraph{Equality transport.}
The major must match $[s,u]$, and opening the proposed formula template
at $u$ must match the target.  The continuation checks its opening at
$s$.  Lemma~\ref{lem:semantic-opening} gives ordinary transport along
the decoded equality.  The argument applies to compound formulas,
internal clauses, and truth aliases as well as equations.

\paragraph{Closure specialization.}
The closure pattern is the triplicate $[t,t,t]$.  The guard forces the
old instance to have numerical image zero; the new instance may have
an arbitrary numerical image.  This is exactly
$B(0)\to\forall yB(y)$.  The zero guard is indispensable.  Without
it, a closure with false zero instance and true successor instances
would turn a proved successor instance into falsity.

\paragraph{Paired elimination.}
The export is always three concatenated copies of the major.  A unary
major $[B]$ therefore exports $[B,B,B]$.  All three copies are opened
identically; the third copy has precisely the witness instance required
by paired elimination.  The newly allocated witness is fresh for the
retained context and conclusion.

An equation major instead exports arity six; a closure major exports
arity nine; a larger array exports still larger arity.  These are all
truth aliases.  Thus failed elimination endpoint claims have only the
effect of adjoining a theorem.  A numeric scalar major, if accepted as
formula padding, exports scalar padding after multiplication and also
denotes truth.  Unsupported types reject.

A ternary formula with unequal annotations still denotes the closure
of its third child.  It fails the triplicate specialization pattern
unless the three fields match.  This loses a possible raw spelling,
not a mathematical soundness condition.  Canonical compilation uses
triplicate closures throughout.

\paragraph{Paired introduction.}
The call \texttt{i(*a,t,2)} uses the function's argument count to enforce
the relevant outer arity.  For a unary goal $a=[B]$, it is exactly
\texttt{i(B,t,2)}: the universal slot is protected while the
existential witness is opened.  For arity at least two, there are too
many arguments and the call rejects.  An empty goal can instead make
the legal two-argument call \texttt{i(t,2)}.  Its original target was
already $\top$, so accepting it cannot add a false theorem.

For the unary goal, the subsequent successor opening gives the intended
premise $B(t,Sz)$.  The local hypothesis is $B(t,0)$ in case mode and
$B(t,z)$ in progress mode.  The freshness invariant supplies the
usual eigenvariable conditions.

\paragraph{Induction remains restricted.}
Among meaningful nonconstant matrix constructors, only binary equality
has even arity.  Unary paired formulas and ternary closures have odd
arity and receive case introduction.  Empty or surplus arrays that
reach the progress branch denote the fixed equation $0=0$, whose
induction is valid.  Scalar padding that cannot provide a length
rejects.  The broader certificate alphabet therefore adds no induction
on a quantified matrix.

\paragraph{Additional fresh parameters.}
An accepted subderivation can contain newly admitted fractional names,
but the older context and target do not depend on them.  Replace these
auxiliary names by zero when closing the rule, retaining its designated
integer eigenvariable where that rule needs one.  The common opening
of an already checked retained goal leaves its decoded meaning
unchanged.  This handles the actual global-floor convention without
pretending that every admissible negative number is an integer.

\subsection{Canonical completeness, including zero witnesses}
\label{subsec:canonical-completeness}

Rename a finite native proof's eigenvariables to fresh negative integers.
Encode its terms with the slot conventions of
Section~\ref{subsec:sorted-arrays}; encode paired formulas as unary and
closures as triplicate.  Hypothesis leaves consume no instruction.
Transport, closure specialization, and paired elimination have their
displayed packets.  Arithmetic computation leaves use
\eqref{eq:base-supply} or \eqref{eq:step-supply} and then close the
target by hypothesis lookup.

One representational issue requires an actual proof transformation.
The packet's truthy witness field selects introduction.  A literal zero
witness therefore cannot simply be inserted into that field.  Define
\[
                 J(t)=I(0,t;v.0).
\]
The computation axiom proves $J(t)=t$, and $J(t)$ has a nonempty array
representation even when its value is zero.  To introduce with a falsy
witness spelling $t$, introduce instead with $J(t)$: transport its
retained assumption to the assumption using $t$, use the original
premise derivation, and transport the resulting successor instance
back to the instance using $J(t)$.  Derived symmetry supplies the
reverse transport.  This is a finite transformation of the premise,
and the paired conclusion itself is unchanged.  If the existential
coordinate is unused, any closed truthy witness works without a
conversion.

The scalar coordinate $1$ is also a semantic zero at numerical root
three, but that observation alone is not a completeness proof.  Closure
specialization and case introduction insert literal raw zero, and
distinct zero aliases need not be interchangeable under a fixed raw
proof spelling.  The computation-backed box gives the required
derivation instead of assuming representational interchangeability.

The mode zero selects progress at an equation.  A different falsy
mode, such as the empty array, selects case.  These are certificate
control tags, not numerical assertions about the witness.  All steps
in the compilation are finite and effective.

\begin{theorem}[Conservativity and completeness of the arity checker]
\label{thm:arity-checker}
Every accepted certificate has its decoded conclusion provable in the
ordinary realization of the native calculus.  For every canonical
native end sequent, the internal checking judgment and the native
calculus have the same consequence relation, in the idealized
unbounded integer-and-array implementation.  In particular, the public
empty-context driver recognizes exactly the canonical closed theorems.
The translation in each direction is effective on finite derivations.
\end{theorem}

\begin{proof}
For soundness, induct on the finite accepting call tree.  Hypothesis
closure uses Lemma~\ref{lem:semantic-opening}; the logical cases are
the audited continuations above.  Arithmetic supply uses
Lemmas~\ref{lem:constructed-supply} and \ref{lem:zero-supply}.
Every discarded annotation is accounted for by its zero or truth
interpretation, and every fresh parameter by the scope invariant.
For completeness, compile a finite native derivation by the preceding
construction, including the boxed-witness transformation.  Its
canonical input expressions require no surplus-arity aliases or fractional
names; generated administrative surplus arrays have exactly the truth
meaning already proved.  This proves both implications without using a finite test
suite as a substitute for induction on derivations.
\end{proof}

The unchanged native calculus is bi-interpretable with full first-order
$\ipra$ by the inherited arithmetic and logical bridge
\cite{generated,corpus}.  In particular,
\[
 (\D_{x,y}B)^\circ
 =\exists x\bigl(B^\circ(x,0)\to\forall yB^\circ(x,y)\bigr).
\]
The whole numerical domain and actual equality are retained.  Its
internally proved numerical round trips and its formula-composite
derivations therefore transfer through
Theorem~\ref{thm:arity-checker}.  We have changed the checked
representation, not replaced bi-interpretability by standard-model
agreement or supplied a new proof of the scalar predecessor theorem.

\subsection{JSON and finite-runtime qualifications}

The primary profile uses an ordinary list-type test.  The smaller byte
profile instead tests multiplication by zero against an empty list.
Among standard JSON runtime values the latter recognizes exactly lists:
numeric values and booleans produce numeric zero, strings produce an
empty string, and null and dictionaries raise rather than becoming
arrays.  This argument assumes ordinary decoded JSON values, not
arbitrary Python objects with overloaded methods.

Finite nonnegative real coordinates have the same interval
interpretation: $[0,3)$ is zero, and $[3+2j,5+2j)$ names slot $j$.
Finite negative coordinates in the admitted range are additional rigid
names.  Equal Python numerical values have the same decoded role,
including Boolean aliases and signed zero.  Nonfinite values admitted
by Python's permissive JSON reader cannot survive a relevant closed
numerical check.  They may be filled as template holes or ignored
inside a constant annotation; neither makes them unchecked parameters.

The alternative opener based on generic iteration has a separate
runtime qualification.  Empty strings and empty dictionaries traverse
as empty arrays; a dictionary whose sole key is empty traverses as a
unary array with an empty child.  Dictionary values are not visited.
Nonempty strings can rely on CPython's recursion limit to reject, as
also occurs in a current corpus comparison profile~\cite{transient}.
The idealized structural-termination claim for the exact-list profile
must not be transferred to that generic profile unchanged.

All executable profiles have ordinary finite memory, stack, and input
limits.  These can reject a sufficiently large otherwise valid proof.
The formal consequence relation is the unbounded finite-certificate
calculus; the delivered program checks certificates within its runtime
resources.  The measurements charge the entire delivered checker,
while the semantic proof and canonical producer are supplied separately.


\clearpage
\begin{thebibliography}{9}
\bibitem{repo} ProveIt contributors, \emph{Concise presentations of
intuitionistic primitive recursive arithmetic}, commit
\texttt{42a48d2f04ad72036e8d7e365ea91ea08c76a9e3}.
\url{https://github.com/VladimirReshetnikov/ProveIt/tree/42a48d2f04ad72036e8d7e365ea91ea08c76a9e3/Logic/PrimitiveRecursiveArithmetic/ConciseIPRA}.
\bibitem{validated} ProveIt incoming archive,
\path{docs/incoming/ipra_validated_core_2026-11-11.zip}, at the commit above.
The exact full native control is copied as
\path{reference/native_typed_full.py}.
\bibitem{erasure} ProveIt inherited scalar-predecessor proof,
\path{Research/certified_zero/reference/scalar_predecessor.tex},
reproduced in Appendix~\ref{sec:hybrid-bootstrap}.
\bibitem{forward} ProveIt inherited equality-reflection proof,
\path{Research/certified_zero/reference/equality_reflection.tex},
reproduced in Appendix~\ref{sec:eic-completion}.
\bibitem{progress} ProveIt inherited primitive-recursion compiler,
\path{Research/certified_zero/reference/primitive_recursion.tex},
reproduced in Appendix~\ref{sec:iterator-coding}.
\bibitem{generated} ProveIt inherited generated-assumption and arity
certificate development, preserved in
\path{Research/certified_zero/reference/inherited_arity_checker.tex}.
\bibitem{constructed} ProveIt research package
\path{Research/constructed_computation}; constructed computation sources
are restated in Appendix~\ref{sec:arity-checker}.
\bibitem{corpus} ProveIt research collection and its package-specific
claim records, all at the pinned commit in reference~\cite{repo}.
\bibitem{transient} ProveIt research package \path{Research/transient_seed},
at the same pinned commit.
\end{thebibliography}
\end{document}
